\documentclass{article}
\usepackage{graphicx} % Required for inserting images

\title{Yes, That's Mine: Asymptotically Foolproof LLM Ownership Identification Against Hidden Adversarial Decoding Parameter Perturbations}
\author{Jerry Bao}
\date{January 2026}
\usepackage[margin=0.75in]{geometry}
\usepackage{abstract}
\usepackage{microtype}
\usepackage{xurl}
\usepackage{hyperref}
\usepackage{mdframed}
\usepackage{amsthm}
\usepackage{amssymb}
\setlength{\parindent}{0pt}
\setlength{\parskip}{0.8em}
\begin{document}

\maketitle

\vspace{-0.5em}
\begin{center}
\small
\textit{Acknowledgments: This research is supported by the Vector Scholarship in Artificial Intelligence, provided through the Vector Institute.}
\end{center}

\begin{abstract}
\small
Creating and training an LLM can take vast amounts of resources, money, and effort. Because of this, many LLM creators seek to protect their credit and IP from being stolen. 'LLM ownership identification', the study of determining whether an anonymous LLM is a stolen version of a known LLM, is therefore one of the primary commercial, legal, and ethical motivators for the research of LLM provenance. White-box methods and backdoor-style black-box methods for LLM ownership identification face challenges with constrained practical assumptions and invasiveness, driving the search for effective and non-invasive black-box methods. There is a practice-theory gap in the current research for non-invasive, black-box methods, and we have proposed a community-wide formal program for classifying threat scenarios which admit 'theoretically foolproof' (in the sense of decidability, statistical consistency, etc.) methods of LLM ownership identification. Starting from first principles, we argued that hidden decoding parameter perturbations are the necessary first model-identity-obfuscation class to investigate for this formal program. We then constructed statistically consistent tests for identifying a black-box LLM's ownership identity under such identity obfuscations. The tests presented in this paper satisfy key desirable properties such as: 1) being non-invasive and fully black box, 2) being statistically consistent, 3) possessing asymptotically fast convergence, and 4) being viable against an adversary with perfect information. Along the way, we implicitly proved a fundamental negative result: standard, token-distribution-modifying decoding parameters are not obfuscating enough to prevent the existence of a non-invasive, theoretically foolproof method for determining a black-box LLM’s ownership identity. In doing so, we established the first constructive evidence of the feasibility of the proposed formal program, which aims to prove similar results under other threat scenarios.
\end{abstract}

\section{Introduction}

Creating and training an LLM can take vast amounts of resources, money, and effort. Because of this, many LLM creators seek to protect their credit and IP from being stolen. 'LLM ownership identification', the study of determining whether an anonymous LLM is a stolen version of a known LLM, is therefore one of the primary commercial, legal, and ethical motivators for the research of LLM provenance.



We briefly survey some of the most common methods in the literature for LLM ownership identification in the \textit{stolen weights }threat scenario, where the goal is figuring out whether LLM B has stolen the weights of LLM A. We will also discuss some current-day challenges each approach faces, providing motivation for the scope of this paper.



\subsection{White-box Methods}

These methods often involve inspecting the parameters of LLM A and LLM B and using a derived signal from said parameters to compare similarity (\url{https://arxiv.org/pdf/2312.04828, https://arxiv.org/pdf/2410.14273, https://arxiv.org/pdf/2511.06390, https://arxiv.org/pdf/2507.03014}). The main limitation of this approach is that they assume white-box access to suspect model LLM B, which is often unrealistic in various practical scenarios, such as when an adversary deliberately hides LLM B's weights or when LLM B is not a stolen model, but its weights are still hidden due privatization and IP protection measures. We will be focusing instead on black-box methods, where it is assumed that the weights of LLM B are not open for inspection.



\subsection{Backdoor-Style Black-Box Methods}

Backdoor-style methods are among the most common black-box ownership identification approaches in the literature (\url{https://arxiv.org/pdf/2401.12255}, \url{https://arxiv.org/pdf/2410.12318}, \url{https://arxiv.org/pdf/2509.03058}, \url{https://arxiv.org/pdf/2407.10887}, \url{https://ojs.aaai.org/index.php/AAAI/article/view/26750}, \url{https://arxiv.org/pdf/2402.14883}, \url{https://arxiv.org/pdf/2509.09703}). They involve training the LLM to behave in an identifiable way (the observable) when a secret input (the trigger) is present in its prompt. Typically, the trigger is a string and the observable is an emitted response, but this is not a universal rule (see \url{https://arxiv.org/pdf/2509.09703} and \url{https://arxiv.org/pdf/2509.03058}).



Backdoor-style methods are considered to be among the most reliable black-box approaches for LLM ownership identification. However, they are not without challenges that may dissuade their usage. We will discuss a few of the practical and theoretical challenges that this approach faces in the current day.



A key concern is that of \textit{harmlessness}, the notion that backdoors should not interfere with general model utility. Backdoor-style methods necessarily require altering model weights and therefore model behavior. Such methods are known as 'invasive' methods. Poorly implemented backdoors risk accidental activation, normal task interference, and general performance degradation. Backdoor method researchers are well-aware of this risk and attempt to maintain model utility, as well as conduct 'Harmlessness' tests to verify that the embedded backdoor does not impact normal task performance. However, harmlessness is primarily empirically, not theoretically validated. As such, they face common challenges with reproducibility and generalizability, even if the initial harmlessness tests show no clear degradation in performance.


As an illustrative example of practical reproducibility challenges, Instructional Fingerprinting (\url{https://arxiv.org/pdf/2401.12255}) reported that their fingerprinting causes no harm in downstream performance based on tests involving multiple models evaluated on multiple benchmarks. However, a more recent paper (\url{https://arxiv.org/pdf/2509.09703}) reported that Instructional Fingerprinting indeed caused notable performance degradations in their tests.



With regards to generalizability, we note that the majority of papers on backdoor-style methods use small to mid-size LLMs. The generalization of harmlessness findings to large models has not thoroughly studied. For example, The Chain and Hash paper (\url{https://arxiv.org/pdf/2407.10887}) observes that the TruthfulQA benchmark (\url{https://arxiv.org/pdf/2109.07958}) demonstrates a significant performance improvement following backdoor embedding compared to other benchmarks. The authors attribute these observed utility gains to the use of diverse prompt formatting templates during the embedding process. It is thus unclear whether these gains mask potential losses and whether they generalize to large, optimized models that have seen many formats in their training already.



Furthermore, we note that harmlessness tests in the literature primarily only focus on performance-based tasks. The interaction and potential interference of these backdoor-style methods with other model objectives such as safety, alignment, and privacy are rarely considered. Systems engineering (Engineering a Safer World: Systems Thinking Applied to Safety) and sociological (Normal Accidents: Living with High-Risk Technologies) paradigms warn against increased operational complexity, noting that complex, not-well-understood interactions of components tend to significantly increase failure risks. We briefly discuss some understudied potential risks here.



Taxonomic studies of jailbreaking (\url{https://arxiv.org/pdf/2308.03825}) and attempts to create universal jailbreaks for LLMs (\url{https://arxiv.org/pdf/2405.20773}) reveal a trend: jailbreaks almost always exploit a contextually-dependent policy that the LLM is intentionally trained on and which is beneficial under normal circumstances (e.g. roleplay framing, developer privileges, etc.). This suggests that introducing contextually-dependent policies switches into models could introduce more potentially exploitable vectors for misalignment, even without any malicious intent to do so. This is of particular relevance to backdoors, which inherently train contextually-dependent policies into the model. Recent work (\url{https://arxiv.org/pdf/2511.12414}) evidences that seemingly benign backdoors can be used as a jailbreak to trigger misaligned behaviors without including any instructions to misalign during the fine-tune. Other studies (\url{https://arxiv.org/pdf/2508.14031}, \url{https://arxiv.org/pdf/2310.03693}, \url{https://arxiv.org/pdf/2502.17424}) have noted unintentional global alignment degradation of models from fine-tunes on seemingly neutral training data. Furthermore, backdoors often require models to conceal their secret behavior from users - the subtle effects of training 'benign deception' into models on downstream model alignment are not well-studied in the context of backdoor-style methods for ownership identification.



Another interference mode that is understudied in this setting is how the backdoors may interact with privacy mechanisms. For example, differential privacy (\url{https://arxiv.org/pdf/1607.00133}) is designed to limit the influence of any single training example on the learned model, which tends to reduce certain forms of memorization. Backdoor methods, on the other hand, rely on memorization-like phenomena, suggesting potential interference risks if not implemented carefully. Studies such as (\url{https://arxiv.org/pdf/2311.06227}, \url{https://arxiv.org/pdf/2411.15831}, \url{https://arxiv.org/pdf/2504.21036}) demonstrate that privacy mechanisms can interfere with backdoors in various situations.


\subsection{Non-Invasive Black-Box Methods}

The importance of cautious investigation of the harmlessness of backdoor-style approaches stems from the fact that they are, fundamentally, \textit{invasive} methods that alter underlying model weights. Ideally, a black-box method of LLM ownership identification is both reliable and non-invasive. This particular paper lies within the domain of \textit{non-invasive, black-box }methods for LLM ownership identification.



There has been a significant body of research in non-invasive, black-box methods for LLM ownership identification. We identify 2 broad flavors of methods: 1) algorithmic/statistical (\url{https://aclanthology.org/2025.findings-acl.546/}, \url{https://arxiv.org/pdf/2510.06605}, \url{https://arxiv.org/pdf/2505.12682}, \url{https://arxiv.org/pdf/2502.00706}, \url{https://arxiv.org/pdf/2405.02466}) and 2) stylistic/behavioral (\url{https://arxiv.org/pdf/2503.01659}, \url{https://arxiv.org/pdf/2408.02871}, \url{https://aclanthology.org/2021.findings-acl.409.pdf}).


There is a conspicuous 'practice-theory gap' in the literature for non-invasive, black-box LLM ownership verification, a situation in which there exist many empirically effective methods but scarce theoretical guarantees for their effectiveness, even within the threat scenarios that they assume. Practice-theory gaps themselves are not uncommon in machine learning, as noted by survey papers such as \url{https://arxiv.org/pdf/1807.03341} and sentiments in the ML community (\url{https://www.science.org/content/article/ai-researchers-allege-machine-learning-alchemy}).



\subsection{A Formal Program}



We attempt to ameliorate this practice-theory gap by proposing a formal program wherein theoreticians attempt to classify the threat scenarios for non-invasive, black-box LLM ownership identification for which an identification method with a 'theoretically foolproof' guarantee (in the sense of decidability, statistical consistency, etc.) exists. To evaluate the preliminary feasibility of such a program, we will attempt to show that such a guarantee exists for a well-motivated, nontrivial threat scenario.



Suppose an adversary were to steal a model's weights. Various perturbations to the deployed model, such as fine-tuning, decoding parameter alterations, system prompt alterations, etc. contribute to obfuscating the model's identity-attributable behavior during black-box tests of provenance. What should be the first type of perturbation class that we investigate? We argue that we should start by studying a \textit{decoding-parameter-only} adversary. We argue this from two lenses: necessity and practicality.



\subsubsection{Necessity}

We identify decoding parameters as being uniquely the only \textit{strictly necessary} part of an LLM's deployment stack other than the LLM's weights that modifies the LLM's statistical behavior. System prompts, safety filters, etc. are not strictly necessary for an LLM to be functional. Indeed, any working LLM must pass the token logit vector through a decoder in order to generate tokens. Thus, any formal analysis of a black-box LLM ownership identification method that claims to admit a theoretically foolproof guarantee must, at some point, confront decoding parameters head-on.



\subsubsection{Practicality}

Decoding parameter perturbations represent one of the most accessible, nontrivial perturbations that obfuscate an LLM's identity-attributable behavior. It is thus one of the most minimal, realistic perturbations that challenge identifiability. Decoding-parameter-only perturbations therefore serve as the opening challenge to the proposed formal program; if this program is to be viable and worthwhile to undertake, we should first demonstrate its ability to withstand this baseline threat scenario.



\subsection{The scope of this paper}

In this paper, we will investigate whether there exists a theoretically foolproof, non-invasive, black-box method for identifying an LLM's ownership identity if we assume that the only allowable model perturbations are hidden perturbations to standard, token-distribution-modifying decoding parameters (defined later in this paper).



We will see that such a concrete method indeed exists. It will furthermore satisfy the following desirable attributes:
\begin{enumerate}
\item Non-invasive and fully black box
\item Statistically consistent
\item Possesses asymptotically fast convergence
\item Works even when the adversary has perfect information
\end{enumerate}



Along the way, we will have proven a fundamental negative result: standard, token-distribution-modifying decoding parameters are not obfuscating enough to prevent the existence of a non-invasive, theoretically foolproof method for determining a black-box LLM’s ownership identity. Furthermore, we will have established the first constructive evidence of the feasibility of the proposed formal program, which aims to prove similar results under other threat scenarios.



\section{Setting}



\subsection{Threat Model: LLM Ownership Identity Obfuscation via Standard Token-Distribution-Modifying Decoding Parameters}

The following formal threat model provides the motivation for the rest of this paper's theoretical analysis:
\begin{itemize}
\item You are the owner of non-reasoning LLM A, which uses a random-sampling-based next-token decoder. You suspect that black-box non-reasoning LLM B, which also uses a random-sampling-based next-token decoder, is a copied instance of LLM A. You do not have access to LLM B's internals (e.g. weights, decoding parameters, etc.). You may call LLM B on whatever prompts you like, however many times as you like.
\item Case 1: LLM B is a copy of LLM A. If LLM B is run by an adversary, the adversary will wish to obfuscate LLM B's provenance and the fact that it is a copy of LLM A by fixing a set of standard token-distribution-modifying decoding parameters one time and concealing their values. We assume the adversary has \textit{perfect information,} and may fix the standard, token-distribution-modifying decoding parameters of LLM B according to this information.
\item Case 2: LLM B is not a copy of LLM A. The actor behind LLM B does not have an adversarial incentive to impersonate LLM A.
\item Excluded case: LLM B is derived from LLM A via a method other than modifying standard token-distribution-modifying decoding parameters.
\item Goal: Determine whether we are in case 1 or 2. I.e. Determine whether or not LLM B is an instance of LLM A "up to decoding parameters".
\end{itemize}



\subsubsection{Justification of assumptions \& Threats not considered in this setting}

Threat vectors not considered in our formal setting include minor model perturbations such as fine-tuning and other methods for deriving models from LLM A other than perturbing standard token-distribution-modifying decoding parameters. The threat model used for our formal analysis explicitly excludes this possibility. Although we exclude fine-tuned or otherwise derived models from the scope of our analysis, we note that the arguments in this paper do not rely on the \textit{full} strength of this assumption. We suggest a strong possibility that our methods may extend to certain carefully-defined subclasses of such models and encourage future work to attempt to cleanly classify practical subclasses of the excluded case for which the arguments and methods in this paper remain theoretically valid.



In addition, we do not consider the case where an adversary may deliberately attempt to have LLM B, when it is not a copy of LLM A, mimic the identity of LLM A. In our threat model, only copied models seek to conceal their identity; models not derived from LLM A are assumed not to engage in adversarial behavioral impersonation. Indeed, positive detection in our setting results in negative consequences for the actor behind LLM B, implying that a rational actor behind LLM B is disincentivized from impersonating LLM A (and it is furthermore in their interests to avoid being flagged as LLM A). This reflects the setting of IP protection rather than tamper or adversarial spoofing detection.



Although not formally analyzed here, we conjecture that, with respect to the identification methods developed in this work, it is extremely hard or infeasible for a model that is not LLM A to reliably impersonate LLM A. A formal treatment of this type of adversary is an interesting direction for future research. To avoid overclaiming, in this paper we restrict mathematical attention to our stated threat model.



\subsection{Refining the threat model assumptions}

\textbf{Refinement 1:} LLM B is assumed to be a non-reasoning LLM whose \textit{only} implementation modifications from the base model are hidden, standard, fixed token-distribution-modifying decoding parameters. In particular, this means that LLM B is not allowed to use non-parametric modifications such as secret system prompts.

\textbf{Refinement 2:} "Standard token-distribution-modifying decoding parameters" is loosely defined as decoding parameters that modify the token distribution without resulting in any deterministic token generations. We list the allowed decoding parameters here:
\begin{itemize}
\item Temperature
\item top-p (nucleus sampling)
\item frequency penalty
\item presence penalty
\item repetition penalty
\item logit\_bias/token\_bias: Let $\displaystyle t$ be a token. Given a token logit $\displaystyle z_{t}$, then a bias $\displaystyle b_{t}$ applied to the token transforms $\displaystyle z_{t}$ to $\displaystyle z_{t} +b_{t}$.
\item response\_format / json\_mode: A constrained-decoding mode that masks all tokens violating a JSON grammar, forcing the output to be valid JSON without modifying remaining logits.
\item tool/function calling toggles: When toggled off, for each token in the tool-calling vocabulary, it transforms its logit $\displaystyle z$ into $\displaystyle z-\lambda $, where $\displaystyle \lambda  >0$ is some constant that can be $\displaystyle \infty $ and is independent of the token.
\item top\_k
\item typical\_p: compute each token's self-information, sort tokens by increasing distance to the distribution's entropy, and keep the smallest prefix of tokens whose probabilities sum to $\geqslant$ typical\_p
\item no\_repeat\_ngram\_size: disallow repeating any n-gram of size no\_repeat\_ngram\_size in the output
\item Mirostat
\item Mirostat 2
\item Top-a/$\displaystyle \eta $-sampling: Truncates tokens that have $\displaystyle p$ values such that $\displaystyle p< p_{max}^{\eta }$
\item min-p: A token logit masking rule
\item tail-free sampling: A token logit masking rule
\end{itemize}



\textbf{TODO:} Add citations for the definitions of each decoding parameter listed above.

\subsection{Definitions}
\begin{itemize}
\item LLM Equality: We say LLM A = LLM B, LLM B is a copy of LLM A, or LLM B is identical to LLM A if LLM A and LLM B share identical architecture, weights, vocabulary, tokenizers, and all other implementation details that constitute a \textit{base} LLM model. This does not include auxiliary deployment layer implementation decisions; in particular, the values of auxiliary deployment parameters such as temperature or logit bias are not included among the implementation details considered when establishing LLM equality.
\item If we have distinct contexts/prompts $\displaystyle C_{1} ,...,C_{m}$, we say that the distinct strings $\displaystyle a_{1} ,...,a_{k} ,b_{1} ,...,b_{k}$ satisfy the $\displaystyle ( \dagger )$ Regularity Conditions with respect to $\displaystyle C_{1} ,...,C_{m}$ if: 1) each of $\displaystyle a_{1} ,...,a_{k} ,b_{1} ,...,b_{k}$ is a token in the vocabulary of LLM A, 2) given any generated output string from LLM A, $\displaystyle t\in \{a_{1} ,...,a_{k} ,b_{1} ,...,b_{k}\}$ is a prefix of that string iff it is the first generated token of that string (practical methods for ensuring this are discussed in the Appendix), 3) none of $\displaystyle a_{1} ,...,a_{k} ,b_{1} ,...,b_{k}$ appears as a token within $\displaystyle C_{j} ,j=1:m$ with respect to LLM A's own tokenization (practical case: none of $\displaystyle a_{1} ,...,a_{k} ,b_{1} ,...,b_{k}$ is a substring of $\displaystyle C_{j} ,j=1:m$), 4) for each $\displaystyle a_{i}$ and each $\displaystyle C_{j}$, it is \textit{possible} for LLM B (after fixing its hidden decoding parameters) to generate an output string on context $\displaystyle C_{j}$ with $\displaystyle a_{i}$ as that output's prefix string, and 5) for each $\displaystyle b_{i}$ and each $\displaystyle C_{j}$, it is \textit{possible} for LLM B (after fixing its hidden decoding parameters) to generate an output string on context $\displaystyle C_{j}$ with $\displaystyle b_{i}$ as that output's prefix string.
\item Suppose we repeatedly call LLM B on each context $\displaystyle C_{j}$. Let $\displaystyle N_{j}$ be the number of times LLM B was called on context $\displaystyle C_{j}$ (determined by us)
\item Let $\displaystyle T$ be LLM B's temperature (unknown to us)
\item For LLM A, let $\displaystyle z_{a_{i} ,j}$ be token $\displaystyle a_{i}$'s base logit given context $\displaystyle C_{j}$. (known to us)
\item For LLM A, Let $\displaystyle z_{b_{i} ,j}$ be token $\displaystyle b_{i}$'s base logit given context $\displaystyle C_{j}$. (known to us)
\item Let $\displaystyle z_{a,j} =\sum _{i=1}^{k} \ z_{a_{i} ,j}$
\item Let $\displaystyle z_{b,j} =\sum _{i=1}^{k} \ z_{b_{i} ,j}$
\item Let $\displaystyle \Delta z_{j} =z_{a,j} -z_{b,j}$
\item Let $\displaystyle p_{j}( a,i)$ be the true probability of $\displaystyle a_{i}$ being a prefix string of a generated output from LLM B on context $\displaystyle C_{j}$ (unknown to us)
\item Let $\displaystyle p_{j}( b,i)$ be the true probability of $\displaystyle b_{i}$ being a prefix string of a generated output from LLM B on context $\displaystyle C_{j}$ (unknown to us)
\item Let $\displaystyle n_{j}( a,i)$ be the actual number of times $\displaystyle a_{i}$ appeared as a prefix string of a generated output from LLM B on context $\displaystyle C_{j}$ (known to us)
\item Let $\displaystyle n_{j}( b,i)$ be the actual number of times $\displaystyle b_{i}$ appeared as a prefix string of a generated output from LLM B on context $\displaystyle C_{j}$ (known to us)
\item Let $\displaystyle g_{j} =ln\left(\prod _{i=1}^{k}\frac{n_{j}( a,i)}{n_{j}( b,i)}\right) =\sum _{i=1}^{k} ln( n_{j}( a,i)) -\sum _{i=1}^{k} ln( n_{j}( b,i))$
\item Let $\displaystyle c_{j} =ln\left(\prod _{i=1}^{k}\frac{p_{j}( a,i)}{p_{j}( b,i)}\right) =\sum _{i=1}^{k} ln( p_{j}( a,i)) -\sum _{i=1}^{k} ln( p_{j}( b,i))$
\end{itemize}



Comment: The $\displaystyle \dagger $ regularity conditions are not intended to be an 'uncontrollable assumption'. They are intended to be \textit{by construction.} The choices of $\displaystyle C_{1} ,...,C_{m}$ and $\displaystyle a_{1} ,...,a_{k} ,b_{1} ,...,b_{k}$ are within the user's control.



We immediately have the following theorem:



\textbf{Theorem 1:} $\displaystyle g_{j}$ is a consistent estimator of $\displaystyle c_{j}$.



Proof: Follows from the fact that $\displaystyle \frac{n_{j}( a,i)}{n_{j}( b,i)} =\frac{n_{j}( a,i) /N_{j}}{n_{j}( b,i) /N_{j}}$ are consistent estimators of $\displaystyle \frac{p_{j}( a,i)}{p_{j}( b,i)}$.\qed 





\section{Theoretical Consequences of LLM A = LLM B}
\begin{mdframed}
Let $C_1, C_2$ be 2 prompts. Let $a_1,\ldots,a_k,b_1,\ldots,b_k$ be $2k$ strings.
For the rest of this section, we shall assume that
$a_1,\ldots,a_k,b_1,\ldots,b_k$ satisfy $\dagger$ with respect to $C_1, C_2$.
\end{mdframed}


Note that by $\displaystyle \dagger ,1$ then $\displaystyle a_{1} ,...,a_{k} ,b_{1} ,...,b_{k}$ simply becomes a list of tokens of LLM B. Furthermore, under $\displaystyle \dagger ,2$, some of our previous definitions equivalently become:
\begin{itemize}
\item Let $\displaystyle p_{j}( a,i)$ be the true probability of $\displaystyle a_{i}$ in the first-token distribution of LLM B with respect to context $\displaystyle C_{j}$ (unknown to us)
\item Let $\displaystyle p_{j}( b,i)$ be the true probability of $\displaystyle b_{i}$ in the first-token distribution of LLM B with respect to context $\displaystyle C_{j}$ (unknown to us)
\item Let $\displaystyle n_{j}( a,i)$ be the actual number of times $\displaystyle a_{i}$ appeared as the first token of LLM B on context $\displaystyle C_{j}$ (known to us)
\item Let $\displaystyle n_{j}( b,i)$ be the actual number of times $\displaystyle b_{i}$ appeared as the first token of LLM B on context $\displaystyle C_{j}$ (known to us)
\end{itemize}



We then provide some additional definitions.

\textbf{Definitions:}
\begin{itemize}
\item Let $\displaystyle b_{a_{i}}$be token $\displaystyle a_{i}$'s logit bias in LLM B (unknown to us)
\item Let $\displaystyle b_{b_{i}}$be token $\displaystyle b_{i}$'s logit bias in LLM B (unknown to us)
\item Let $\displaystyle b_{a} =\sum _{i=1}^{k} b_{a_{i}}$
\item Let $\displaystyle b_{b} =\sum _{i=1}^{k} b_{b_{i}}$
\item Let $\displaystyle \Delta b=b_{a} -b_{b}$
\end{itemize}



\textbf{Theorem 2:} Suppose LLM A = LLM B. Fix the context to be $\displaystyle C_{1}$ or $\displaystyle C_{2}$. For all $\displaystyle i$ in $\displaystyle 1:k$, the value $\displaystyle \frac{p_{j}( a,i)}{p_{j}( b,i)}$ is invariant under all standard token-distribution-modifying decoding parameters applied to LLM B, except for temperature $\displaystyle T$ and logit biases $\displaystyle b_{a_{1}} ,...,b_{a_{k}} ,b_{b_{1}} ...,b_{b_{k}}$.



Proof:

We first consider the following decoding parameters:
\begin{itemize}
\item top-p (nucleus sampling)
\item response\_format / json\_mode: A constrained-decoding mode that masks all tokens violating a JSON grammar, forcing the output to be valid JSON without modifying remaining logits.
\item top\_k
\item typical\_p: compute each token's self-information, sort tokens by increasing distance to the distribution's entropy, and keep the smallest prefix of tokens whose probabilities sum to $\displaystyle \geqslant $typical\_p
\item no\_repeat\_ngram\_size: disallow repeating any n-gram of size no\_repeat\_ngram\_size in the output
\item Mirostat
\item Top-a/$\displaystyle \eta $-sampling: Truncates tokens that have $\displaystyle p$ values such that $\displaystyle p< p_{max}^{\eta }$
\item min-p: A token logit masking rule
\item tail-free sampling: A token logit masking rule
\end{itemize}

These are all token masking decoding parameters. In other words, they each define a set $\displaystyle S$ of allowed tokens. Note that $\displaystyle a_{i}$ and $\displaystyle b_{i}$ are assumed to be in $\displaystyle S$ due to $\displaystyle \dagger ,4$ and $\displaystyle \dagger ,5$. 

Before masking, $\displaystyle \frac{p_{j}( a,i)}{p_{j}( b,i)} =\frac{\frac{e^{z_{a_{i}}}}{\sum _{t} e^{z_{t}}}}{\frac{e^{z_{b_{i}}}}{\sum _{t} e^{z_{t}}}} =\frac{e^{z_{a_{i}}}}{e^{z_{b_{i}}}}$, where $\displaystyle z$ is each token's logit. After masking, $\displaystyle \frac{p_{j}( a,i)}{p_{j}( b,i)} =\frac{\frac{e^{z_{a_{i}}}}{\sum _{t\in S} e^{z_{t}}}}{\frac{e^{z_{b_{i}}}}{\sum _{t\in S} e^{z_{t}}}} =\frac{e^{z_{a_{i}}}}{e^{z_{b_{i}}}}$. We conclude that none of the masking decoding parameters affects $\displaystyle \frac{p_{j}\left( a,i\right)}{p_{j}\left( b,i\right)}$.



We next consider the following 3 decoding parameters:
\begin{itemize}
\item frequency penalty
\item presence penalty
\item repetition penalty
\end{itemize}

Note that these 3 decoding parameters only kick in if $\displaystyle a_{i}$ or $\displaystyle b_{i}$ actually appear in the context. But by $\displaystyle \dagger ,3$, neither of them appears in the context. Therefore, neither $\displaystyle z_{a_{i}}$ or $\displaystyle z_{b_{i}}$ is modified. Using a similar argument as before, we see that $\displaystyle \frac{p_{j}\left( a,i\right)}{p_{j}\left( b,i\right)} =\frac{e^{z_{a_{i}}}}{e^{z_{b_{i}}}}$ both before and after each of these decoding parameters is applied. We conclude that none of these 3 decoding parameters affects $\displaystyle \frac{p_{j}\left( a,i\right)}{p_{j}\left( b,i\right)}$.



We next consider the following decoding parameter:
\begin{itemize}
\item tool/function calling toggles: When toggled off, for each token in the tool-calling vocabulary, it transforms its logit $\displaystyle z$ into $\displaystyle z-\lambda $, where $\displaystyle \lambda  >0$ is some constant that can be $\displaystyle \infty $ and is independent of the token.
\end{itemize}

Note that if $\displaystyle \lambda =\infty $, then this is a masking decoding parameter. In this case, the same reasoning as before applies for why $\displaystyle \frac{p_{j}\left( a,i\right)}{p_{j}\left( b,i\right)}$ is unaffected. Else, this is just a standard token logit bias, which is permitted within the assumptions of theorem 2.



Finally, we consider the following decoding parameter.
\begin{itemize}
\item Mirostat 2
\end{itemize}

Note that Mirostat 2 does not affect the logit values of surviving tokens in the first-token distribution. This is because Mirostat 2 needs adaptive feedback from earlier tokens in the output and none exist for the first token. We can therefore apply our reasoning from the masking case to conclude that $\displaystyle \frac{p_{j}\left( a,i\right)}{p_{j}\left( b,i\right)}$ is unaffected by Mirostat 2.



Finally, we can conclude that only temperature and logit bias affect the ratio $\displaystyle \frac{p_{j}\left( a,i\right)}{p_{j}\left( b,i\right)}$. This completes the proof of the claim.\qed 



\textbf{Corollary 1:} Suppose LLM A = LLM B. Fix the context to be $\displaystyle C_{1}$ or $\displaystyle C_{2}$. The value $\displaystyle \prod _{i=1}^{k}\frac{p_{j}\left( a,i\right)}{p_{j}\left( b,i\right)}$ (and therefore $\displaystyle c_{j}$) with respect to the first-token distribution is invariant under all standard token-distribution-modifying decoding parameters, except for temperature and bias.\qed 



\textbf{Corollary 2:} Suppose LLM A = LLM B. Then $\displaystyle c_{j} =ln\left(\prod _{i=1}^{k}\frac{p_{j}\left( a,i\right)}{p_{j}\left( b,i\right)}\right) =\frac{\Delta z_{j} +\Delta b}{T}$

Proof:

Since only temperature and logit bias affect probability ratios (theorem 2), when computing $\displaystyle \frac{p_{j}\left( a,i\right)}{p_{j}\left( b,i\right)}$, we may assume only only temperature and bias affect the token logits. Thus:

$\displaystyle \begin{array}{l}
\frac{p_{j}\left( a,i\right)}{p_{j}\left( b,i\right)} =\frac{exp\left(\frac{z_{a_{i} ,j} +b_{a_{i}}}{T}\right) /\sum _{t} exp\left(\frac{z_{t,j} +b_{t}}{T}\right)}{exp\left(\frac{z_{b_{i} ,j} +b_{b_{i}}}{T}\right) /\sum _{t} exp\left(\frac{z_{t,j} +b_{t}}{T}\right)}\\
=exp\left(\frac{\left( z_{a_{i} ,j} -z_{b_{i} ,j}\right) +\left( b_{a_{i}} -b_{b_{i}}\right)}{T}\right)
\end{array}$

Thus:

$\displaystyle  \begin{array}{l}
\prod _{i=1}^{k}\frac{p_{j}\left( a,i\right)}{p_{j}\left( b,i\right)} =\prod _{i=1}^{k} exp\left(\frac{\left( z_{a_{i} ,j} -z_{b_{i} ,j}\right) +\left( b_{a_{i}} -b_{b_{i}}\right)}{T}\right)\\
=exp\left(\sum _{i=1}^{k}\frac{\left( z_{a_{i} ,j} -z_{b_{i} ,j}\right) +\left( b_{a_{i}} -b_{b_{i}}\right)}{T}\right)\\
=exp\left(\frac{\Delta z_{j} +\Delta b}{T}\right)
\end{array}$

Therefore, $\displaystyle c_{j} =ln\left(\prod _{i=1}^{k}\frac{p_{j}\left( a,i\right)}{p_{j}\left( b,i\right)}\right) =\frac{\Delta z_{j} +\Delta b}{T}$.\qed 



\subsection{Temperature Estimation}

\textbf{Theorem 3:} Suppose LLM A = LLM B. Then $\displaystyle T=\frac{\Delta z_{1} -\Delta z_{2}}{c_{1} -c_{2}}$

Proof:

$\displaystyle T=\frac{\Delta z_{1} -\Delta z_{2}}{c_{1} -c_{2}}$

$\displaystyle =\frac{\Delta z_{1} -\Delta z_{2}}{\frac{\Delta z_{1} +\Delta b}{T} -\frac{\Delta z_{2} +\Delta b}{T}}$ (by corollary 2)

$\displaystyle =T\cdot \frac{\Delta z_{1} -\Delta z_{2}}{\Delta z_{1} -\Delta z_{2}}$ (1)

$\displaystyle =T$\qed 





\textbf{Theorem 4:} Suppose LLM A = LLM B. Then $\displaystyle \frac{\Delta z_{1} -\Delta z_{2}}{g_{1} -g_{2}}$ is a consistent estimator of $\displaystyle T$.

Proof:

This follows from theorem 1 and theorem 3.\qed 





\subsection{Logit Bias Displacement Estimation}

\textbf{Theorem 5:} Suppose LLM A = LLM B. Then $\displaystyle \Delta b=\frac{c_{2} \Delta z_{1} -c_{1} \Delta z_{2}}{c_{1} -c_{2}}$

Proof:

By Corollary 2, we have:

$\displaystyle  \begin{array}{l}
\frac{c_{2} \Delta z_{1} -c_{1} \Delta z_{2}}{c_{1} -c_{2}} =\frac{\frac{\Delta z_{2} +\Delta b}{T} \Delta z_{1} -\frac{\Delta z_{1} +\Delta b}{T} \Delta z_{2}}{\frac{\Delta z_{1} +\Delta b}{T} -\frac{\Delta z_{2} +\Delta b}{T}}\\
=\frac{\left( \Delta z_{2} +\Delta b\right) \Delta z_{1} -\left( \Delta z_{1} +\Delta b\right) \Delta z_{2}}{\Delta z_{1} -\Delta z_{2}}\\
=\Delta b
\end{array}$



\qed 

\textbf{Theorem 6:} Suppose LLM A = LLM B. Then $\displaystyle \frac{g_{2} \Delta z_{1} -g_{1} \Delta z_{2}}{g_{1} -g_{2}}$ is a consistent estimator of $\displaystyle \Delta b$.

Proof: This follows from theorem 1 and theorem 5.\qed 



\subsection{Additional remarks for this section}

Let us remind ourselves what $\displaystyle \Delta b$ is: it is the difference between the sum of the logit biases of the $\displaystyle a_{i}$'s and the sum of the logit biases of the $\displaystyle b_{i}$'s. If $\displaystyle k=1$, then $\displaystyle \Delta b$ is simply the difference between the logit biases of $\displaystyle a_{1}$ and $\displaystyle b_{1}$. Let $\displaystyle a:=a_{1}$ and $\displaystyle b:=b_{1}$. In this case, we may write $\displaystyle \Delta b$ as $\displaystyle \Delta b_{a-b}$. Introduce a 3rd token $\displaystyle c$. Using theorem 6, we can asymptotically compute $\displaystyle \Delta b_{a-b} ,\Delta b_{b-c} ,\Delta b_{c-a}$. This gives us a system of 3 linear equations where the variables are the logit biases $\displaystyle b_{a}$, $\displaystyle b_{b}$, and $\displaystyle b_{c}$. Solving gives us the logit biases of $\displaystyle a,b,c$.



Our discussion informs us that, assuming LLM A = LLM B, we can theoretically recover the individual logit biases that were applied to specific tokens in LLM B to an arbitrarily high degree of accuracy. This is perhaps a somewhat surprising identifiability result. Moreover, combining this with theorem 4 tells us that when LLM A = LLM B, logit biases and temperature are both recoverable information, even if we did not initially know which standard token-distribution-modifying decoding parameters LLM B used.



\section{Theoretical Consequences of LLM A $\neq$ LLM B}
\begin{mdframed}
Let $C_1, C_2$ be 2 prompts. Let $a_1,\ldots,a_k,b_1,\ldots,b_k$ be $2k$ strings.
For the rest of this section, we shall assume that
$a_1,\ldots,a_k,b_1,\ldots,b_k$ satisfy $\dagger$ with respect to $C_1, C_2$.
\end{mdframed}


Definitions:
\begin{itemize}
\item Let $\displaystyle est\left( C_{1} ;C_{2} ;a_{1} ,...,a_{k} ;b_{1} ,...,b_{k}\right) =\frac{\Delta z_{1} -\Delta z_{2}}{g_{1} -g_{2}}$
\item Let $\displaystyle est2\left( C_{1} ;C_{2} ;a_{1} ,...,a_{k} ;b_{1} ,...,b_{k}\right) =\frac{g_{2} \Delta z_{1} -g_{1} \Delta z_{2}}{g_{1} -g_{2}}$
\end{itemize}

It is trivial to see that even though LLM A $\displaystyle \neq $ LLM B, $\displaystyle est$ and $\displaystyle est2$ are still consistent estimators. Namely, we have:



Corollary 3: $\displaystyle est\left( C_{1} ;C_{2} ;a_{1} ,...,a_{k} ;b_{1} ,...,b_{k}\right)$ is a consistent estimator of $\displaystyle \frac{\Delta z_{1} -\Delta z_{2}}{c_{1} -c_{2}}$.\qed 



Corollary 4: $\displaystyle est2\left( C_{1} ;C_{2} ;a_{1} ,...,a_{k} ;b_{1} ,...,b_{k}\right)$ is a consistent estimator of $\displaystyle \frac{c_{2} \Delta z_{1} -c_{1} \Delta z_{2}}{c_{1} -c_{2}}$.\qed 



\section{Determining LLM Ownership Identity}

We now have the tools to construct specific tests for LLM ownership identity. We will first present the most general form of the framework, then illustrate specific examples afterwards.



But first, let us introduce some new definitions that generalize our old definitions:



\textbf{Definitions:}
\begin{itemize}
\item If we have contexts/prompts $\displaystyle C_{1} ,...,C_{m}$, we say that the strings $\displaystyle a_{1,j} ,...,a_{k_{j} ,j} ,b_{1,j} ,...,b_{k_{j} ,j}$ for $\displaystyle j=1:m$ satisfy the $\displaystyle \left( \dagger \dagger \right)$ Regularity Conditions with respect to $\displaystyle C_{1} ,...,C_{m}$ if: 1) for each $\displaystyle a_{i,j}$ and each $\displaystyle C_{j}$, it is \textit{possible} for LLM B (after fixing its hidden decoding parameters) to generate an output string on context $\displaystyle C_{j}$ with $\displaystyle a_{i,j}$ as that output's prefix string, and 2) for each $\displaystyle b_{i,j}$ and each $\displaystyle C_{j}$, it is \textit{possible} for LLM B (after fixing its hidden decoding parameters) to generate an output string on context $\displaystyle C_{j}$ with $\displaystyle b_{i,j}$ as that output's prefix string.
\item Suppose we repeatedly call LLM B on each context $\displaystyle C_{j}$. Let $\displaystyle N_{j}$ be the number of times LLM B was called on context $\displaystyle C_{j}$ (determined by us).
\item Let $\displaystyle p_{j}\left( a,i\right)$ be the true probability of $\displaystyle a_{i,j}$ being a prefix string of a generated output from LLM B on context $\displaystyle C_{j}$ (unknown to us)
\item Let $\displaystyle p_{j}\left( b,i\right)$ be the true probability of $\displaystyle b_{i,j}$ being a prefix string of a generated output from LLM B on context $\displaystyle C_{j}$ (unknown to us)
\item Let $\displaystyle n_{j}\left( a,i\right)$ be the actual number of times $\displaystyle a_{i,j}$ appeared as a prefix string of a generated output from LLM B on context $\displaystyle C_{j}$ (known to us)
\item Let $\displaystyle n_{j}\left( b,i\right)$ be the actual number of times $\displaystyle b_{i,j}$ appeared as a prefix string of a generated output from LLM B on context $\displaystyle C_{j}$ (known to us)
\item Let $\displaystyle g_{j} =ln\left(\prod _{i=1}^{k_{j}}\frac{n_{j}\left( a,i\right)}{n_{j}\left( b,i\right)}\right) =\sum _{i=1}^{k_{j}} ln\left( n_{j}\left( a,i\right)\right) -\sum _{i=1}^{k_{j}} ln\left( n_{j}\left( b,i\right)\right)$
\item Let $\displaystyle c_{j} =ln\left(\prod _{i=1}^{k_{j}}\frac{p_{j}\left( a,i\right)}{p_{j}\left( b,i\right)}\right) =\sum _{i=1}^{k_{j}} ln\left( p_{j}\left( a,i\right)\right) -\sum _{i=1}^{k_{j}} ln\left( p_{j}\left( b,i\right)\right)$
\end{itemize}



Comment: The $\displaystyle \dagger \dagger $ regularity conditions are not intended to be an 'uncontrollable assumption'. They are intended to be \textit{by construction.} The choices of $\displaystyle C_{1} ,...,C_{m}$ and $\displaystyle a_{1,j} ,...,a_{k_{j} ,j} ,b_{1,j} ,...,b_{k_{j} ,j} ;j=1:m$ are within the user's control.



We must also introduce the concept of an \textit{experiment.}



\subsection{What is an experiment?}

Note that each $\displaystyle n_{j}\left( a,i\right)$ comes from running an experiment that calls LLM B multiple times on context $\displaystyle C_{j}$ and recording the number of times that $\displaystyle a_{i,j}$ appears as an output prefix. Similar story for $\displaystyle n_{j}\left( b,i\right)$. It must be clarified that the actual sequence of runs we make with LLM B on context $\displaystyle C_{j}$ might be \textit{shared} between different $\displaystyle j$'s.



Consider the following illustrative scenario: Fix context $\displaystyle C$ and let $\displaystyle C_{j_{1}} =C_{j_{2}} =C$. For $\displaystyle C_{j_{1}}$, associate the strings $\displaystyle a_{1,j_{1}} ,...,a_{k_{j_{1}} ,j_{1}} ,b_{1,j_{1}} ,...,b_{k_{j_{1}} ,j_{1}}$. For $\displaystyle C_{j_{2}}$, associate the strings $\displaystyle a_{1,j_{2}} ,...,a_{k_{j_{2}} ,j_{2}} ,b_{1,j_{2}} ,...,b_{k_{j_{2}} ,j_{2}}$, which may or may not overlap with $\displaystyle a_{1,j_{1}} ,...,a_{k_{j_{1}} ,j_{1}} ,b_{1,j_{1}} ,...,b_{k_{j_{1}} ,j_{1}}$. Now run LLM B repeatedly on context $\displaystyle C=C_{j_{1}} =C_{j_{2}}$. This sequence of runs constitutes a single experiment. From this single experiment, we may record the number of output prefix occurrences of $\displaystyle a_{1,j_{1}} ,...,a_{k_{j_{1}} ,j_{1}} ,b_{1,j_{1}} ,...,b_{k_{j_{1}} ,j_{1}}$ (i.e. record $\displaystyle n_{j_{1}}\left( a,i\right)$ and $\displaystyle n_{j_{1}}\left( b,i\right)$) and the number of output prefix occurrences of $\displaystyle \ a_{1,j_{2}} ,...,a_{k_{j_{2}} ,j_{2}} ,b_{1,j_{2}} ,...,b_{k_{j_{2}} ,j_{2}}$ (i.e. record $\displaystyle n_{j_{2}}\left( a,i\right)$ and $\displaystyle n_{j_{2}}\left( b,i\right)$). Crucially, note that $\displaystyle g_{j_{1}}$ and $\displaystyle g_{j_{2}}$ are not built from independent, self-contained experiments. They are built from the \textit{same} experiment and therefore coupled and non-independent.



The key point is this: different $\displaystyle j$'s do not necessarily represent different experiments. Different $\displaystyle j$'s may represent recorded information from the \textit{same} experiment. Not merely from a\textit{ copy} of an experiment, but from the \textit{same} experiment.



We can now define what an experiment is.



\textbf{Definitions:}
\begin{itemize}
\item An experiment $\displaystyle E$ encapsulates all the information about the full sequence of runs we do with LLM B on a fixed context. Output prefix count data is recorded from this experiment and may be used for multiple different $\displaystyle j$'s.
\item Let $\displaystyle E_{1} ,...,E_{w}$ be the list of experiments we have done.
\item For each $\displaystyle r=1:w$, let $\displaystyle O_{r}$ be the subset of $\displaystyle n_{j}\left( a,i\right)$ and $\displaystyle n_{j}\left( b,i\right)$ observations that came from experiment $\displaystyle E_{r}$.
\item Let $\displaystyle k^{\left( r\right)}$ be the number of elements in $\displaystyle O^{\left( r\right)}$
\end{itemize}



Here are some easy facts about $\displaystyle O_{r}$:
\begin{itemize}
\item The $\displaystyle O_{r}$'s partition the $\displaystyle n_{j}\left( a,i\right)$'s and $\displaystyle n_{j}\left( b,i\right)$'s. After all, each $\displaystyle n_{j}\left( a,i\right)$ or $\displaystyle n_{j}\left( b,i\right)$ may only come from a single experiment $\displaystyle E_{r}$.
\item The different $\displaystyle O_{r}$'s are statistically independent of each other. After all, they record observations from independent experiments.
\item For each $\displaystyle j$, $\displaystyle n_{j}\left( a,i\right)$ and $\displaystyle n_{j}\left( b,i\right)$ for $\displaystyle i=1:k_{j}$ all belong to the same $\displaystyle O_{r}$. This justifies the notation $\displaystyle j\in E_{r}$, which is taken to mean that the observations for $\displaystyle j$ (all) came from $\displaystyle E_{r}$.
\item Each $\displaystyle O_{r}$ is associated with a number $\displaystyle N^{\left( r\right)}$ representing the number of times LLM B was run/called during the experiment. Note that for any $\displaystyle j=1:m$, if $\displaystyle j\in E_{r}$, then $\displaystyle N^{\left( r\right)} =N_{j}$.
\item For any $\displaystyle j_{1} ,j_{2}$, if $\displaystyle j_{1} \in E_{r}$ and $\displaystyle j_{2} \in E_{r}$, then $\displaystyle C_{j_{1}} =C_{j_{2}}$.
\item For any $\displaystyle j=1:m$, $\displaystyle j\in E_{r}$ for some $\displaystyle r=1:w$.
\end{itemize}



We now have the language to talk about the general LLM ownership identification method.



\subsection{The General Method}
\begin{mdframed}
Throughout this "The General Method" subsection, we will be assuming $\dagger\dagger$ instead of $\dagger$ (which is a special case of $\dagger\dagger$; $\dagger\dagger$ is a relaxation of $\dagger$).
\end{mdframed}

\subsubsection{The Estimator Algorithm}
\begin{enumerate}
\item For each $\displaystyle j$, construct estimator $\displaystyle g_{j} =ln\left(\prod _{i=1}^{k_{j}}\frac{n_{j}\left( a,i\right)}{n_{j}\left( b,i\right)}\right)$
\item Construct a real-analytic function $\displaystyle \phi \left( x_{1} ,...,x_{m}\right) :\ \mathbb{R}^{m}\rightarrow \mathbb{R}$ such that $\displaystyle \phi \left( g_{1} ,...,g_{m}\right)$ is a consistent estimator of a known constant $\displaystyle C$ if LLM A = LLM B. $\displaystyle \phi \left( x_{1} ,...,x_{m}\right)$ must not be identically $\displaystyle C$.
\item \textbf{Decision rule:} If $\displaystyle \phi \left( g_{1} ,...,g_{m}\right)$ converges to $\displaystyle C$ as $\displaystyle N^{\left( 1\right)} ,...,N^{\left( w\right)}\rightarrow \infty $, conclude that LLM A = LLM B. Else, conclude that LLM A $\displaystyle \neq $LLM B.
\end{enumerate}



\textbf{Properties of the estimator algorithm:}
\begin{enumerate}
\item  $\displaystyle \phi \left( g_{1} ,...,g_{m}\right)$ is a convergent statistic as $\displaystyle N^{\left( 1\right)} ,...,N^{\left( w\right)}\rightarrow \infty $.
\item If LLM A = LLM B, this statistic converges to $\displaystyle C$. Else, it converges to a non-$\displaystyle C$ value with probability 1.
\end{enumerate}

We shall defer the proof of the estimator algorithm's properties to the end of this section $\displaystyle \left( \star \right)$.



We can also use this same $\displaystyle \phi \left( g_{1} ,...,g_{m}\right)$ to construct a hypothesis test:



\subsubsection{The Hypothesis Test}

Let $\displaystyle \phi \left( g_{1} ,...,g_{m}\right)$ be from step 2 of the estimator algorithm.

Pick any function $\displaystyle \epsilon \left( N^{\left( 1\right)} ,...,N^{\left( w\right)}\right)$ such that $\displaystyle \epsilon \downarrow 0$ and $\displaystyle \frac{\epsilon }{\sqrt{\sum _{r=1}^{w} 1/N^{\left( r\right)}}}\rightarrow \infty $ as $\displaystyle \ N^{\left( 1\right)} ,...,N^{\left( w\right)}\rightarrow \infty $ (e.g. $\displaystyle \epsilon \left( N^{\left( 1\right)} ,...,N^{\left( w\right)}\right) =min\left( N^{\left( 1\right)} ,...,N^{\left( w\right)}\right)^{-\alpha }$ where $\displaystyle 0< \alpha < 1/2$)

Let $\displaystyle H_{0}$ be the null hypothesis that LLM A = LLM B.

Let $\displaystyle H_{1}$ be the alternate hypothesis that LLM A$\displaystyle \neq $LLM B.

Assume $\displaystyle N^{\left( 1\right)} ,...,N^{\left( w\right)}\rightarrow \infty $.



\textbf{Decision rule:} If $\displaystyle |\phi \left( g_{1} ,...,g_{m}\right) -C| >\epsilon \left( N^{\left( 1\right)} ,...,N^{\left( w\right)}\right)$, reject $\displaystyle H_{0}$. Else, fail to reject $\displaystyle H_{0}$.



\textbf{Properties of the hypothesis test:}
\begin{enumerate}
\item If $\displaystyle H_{0}$ if true, with probability 1, the test does not reject $\displaystyle H_{0}$ as $\displaystyle N^{\left( 1\right)} ,...,N^{\left( w\right)}\rightarrow \infty $.
\item If $\displaystyle H_{1}$ is true, with probability 1, the test rejects $\displaystyle H_{0}$ as $\displaystyle N^{\left( 1\right)} ,...,N^{\left( w\right)}\rightarrow \infty $.
\end{enumerate}

Thus, the hypothesis test is consistent in the sense of properties 1 and 2. We shall defer the proof of the hypothesis test's properties to the end of this section $\displaystyle \left( \star \right)$.



\subsection{Specific Examples}

We will present 2 illustrative methods for determining LLM ownership identity using the general framework that was just presented. In particular, we will show that it is possible to construct the estimator $\displaystyle \phi \left( g_{1} ,...,g_{m}\right)$ and corresponding known constant $\displaystyle C$ which form the crux of the estimator algorithm and hypothesis test.



The first estimator will be built from $\displaystyle est$ and the second estimator will be built from $\displaystyle est2$. Intuitively, the $\displaystyle \phi \left( g_{1} ,...,g_{m}\right)$ estimator that we construct will exploit a sharp dichotomy between the cases of LLM A = LLM B and LLM A$\displaystyle \neq $LLM B: conditional stability. When LLM A = LLM B, then $\displaystyle est$ and $\displaystyle est2$ will always converge to fixed and stable quantities $\displaystyle T$ and $\displaystyle \Delta b$. However, when LLM A$\displaystyle \neq $LLM B, the limits of $\displaystyle est$ and $\displaystyle est2$ will vary noisily and unpredictably, determined by how the specific instances of $\displaystyle est$ and $\displaystyle est2$ were constructed.



Because the goal here is to be illustrative and to demonstrate the \textit{existence }of such methods, do not expect the estimators constructed in this section to be the fastest-converging estimators that can be constructed from $\displaystyle est$ and $\displaystyle est2$.



\subsubsection{Naive Method I}

Step 1:

Take contexts $\displaystyle C_{1}$, $\displaystyle C_{2}$ and strings $\displaystyle a_{1} ,...,a_{k}$ and $\displaystyle b_{1} ,...,b_{k}$ satisfying the ($\displaystyle \dagger $) conditions. Similarly, take contexts $\displaystyle C_{3} ,C_{4}$ and strings $\displaystyle a_{1} ',...,a_{k'} '$ and $\displaystyle b_{1} ',...,b_{k'} '$ satisfying the ($\displaystyle \dagger $) conditions.



Step 2:

Relabel to the notation of this section:
\begin{itemize}
\item Let $\displaystyle \langle a_{1,1} ,...,a_{k_{1} ,1} ,b_{1,1} ,...,b_{k_{1} ,1} \rangle =\langle a_{1,2} ,...,a_{k_{2} ,2} ,b_{1,2} ,...,b_{k_{2} ,2} \rangle :=\langle a_{1} ,...,a_{k} ,b_{1} ,...,b_{k} \rangle $
\item Let $\displaystyle \langle a_{1,3} ,...,a_{k_{3} ,3} ,b_{1,3} ,...,b_{k_{3} ,3} \rangle =\langle a_{1,4} ,...,a_{k_{4} ,4} ,b_{1,4} ,...,b_{k_{4} ,4} \rangle :=\langle a_{1} ',...,a_{k'} ',b_{1} ',...,b_{k'} '\rangle $
\end{itemize}



Step 3:

For each $\displaystyle j=1:4$, construct estimator $\displaystyle g_{j} =ln\left(\prod _{i=1}^{k_{j}}\frac{n_{j}\left( a,i\right)}{n_{j}\left( b,i\right)}\right)$ using $\displaystyle C_{j}$ and $\displaystyle a_{1,j} ,...,a_{k_{j} ,j} ,b_{1,j} ,...,b_{k_{j} ,j}$.



Step 4:

Define:
\begin{itemize}
\item For LLM A, let $\displaystyle z_{a_{i,j}}$ be token $\displaystyle a_{i,j}$'s base logit given context $\displaystyle C_{j}$. (known to us)
\item For LLM A, Let $\displaystyle z_{b_{i,j}}$ be token $\displaystyle b_{i,j}$'s base logit given context $\displaystyle C_{j}$. (known to us)
\item Let $\displaystyle z_{a,j} =\sum _{i=1}^{k} \ z_{a_{i,j}}$
\item Let $\displaystyle z_{b,j} =\sum _{i=1}^{k} \ z_{b_{i,j}}$
\item Let $\displaystyle \Delta z_{j} =z_{a,j} -z_{b,j}$
\end{itemize}

Choose $\displaystyle \phi \left( x_{1} ,x_{2} ,x_{3} ,x_{4}\right) =\frac{\Delta z_{1} -\Delta z_{2}}{x_{1} -x_{2}} -\frac{\Delta z_{3} -\Delta z_{4}}{x_{3} -x_{4}}$.



Step 5:

We have $\displaystyle \phi \left( g_{1} ,g_{2} ,g_{3} ,g_{4}\right) =est\left( C_{1} ;C_{2} ;a_{1} ,...,a_{k} ;b_{1} ,...,b_{k}\right) -est\left( C_{3} ;C_{4} ;a_{1} ',...,a_{k'} ';b_{1} ',...,b_{k'} '\right)$.

Let $\displaystyle C=0$.



\textbf{Claim:} $\displaystyle \phi \left( g_{1} ,g_{2} ,g_{3} ,g_{4}\right)$ is a consistent estimator of $\displaystyle C=0$ if LLM A = LLM B.



Proof:

Assume LLM A = LLM B.

By theorem 4, both $\displaystyle est\left( C_{1} ;C_{2} ;a_{1} ,...,a_{k} ;b_{1} ,...,b_{k}\right)$ and $\displaystyle est\left( C_{3} ;C_{4} ;a_{1} ',...,a_{k'} ';b_{1} ',...,b_{k'} '\right)$ converge to $\displaystyle T$, so $\displaystyle \phi \left( g_{1} ,g_{2} ,g_{3} ,g_{4}\right)$ converges to 0.\qed 



\subsubsection{Naive Method II}

Step 1:

Take contexts $\displaystyle C_{1} ,C_{2} ,C_{3} ,C_{4}$ and strings $\displaystyle a_{1} ,...,a_{k}$ and $\displaystyle b_{1} ,...,b_{k}$ satisfying the $\displaystyle \dagger $ regularity conditions for both $\displaystyle \left( C_{1} ,C_{2}\right)$ and $\displaystyle \left( C_{3} ,C_{4}\right)$.



Step 2:

Relabel to the notation of this section:

$\displaystyle \langle a_{1,j} ,...,a_{k_{j} ,j} ,b_{1,j} ,...,b_{k_{j} ,j} \rangle :=\langle a_{1} ,...,a_{k} ,b_{1} ,...,b_{k} \rangle $ for $\displaystyle j=1:4$.



Step 3:

For each $\displaystyle j=1:4$, construct estimator $\displaystyle g_{j} =ln\left(\prod _{i=1}^{k_{j}}\frac{n_{j}\left( a,i\right)}{n_{j}\left( b,i\right)}\right)$ using $\displaystyle C_{j}$ and $\displaystyle a_{1,j} ,...,a_{k_{j} ,j} ,b_{1,j} ,...,b_{k_{j} ,j}$.



Step 4:

Define:
\begin{itemize}
\item For LLM A, let $\displaystyle z_{a_{i,j}}$ be token $\displaystyle a_{i,j}$'s base logit given context $\displaystyle C_{j}$. (known to us)
\item For LLM A, Let $\displaystyle z_{b_{i,j}}$ be token $\displaystyle b_{i,j}$'s base logit given context $\displaystyle C_{j}$. (known to us)
\item Let $\displaystyle z_{a,j} =\sum _{i=1}^{k} \ z_{a_{i,j}}$
\item Let $\displaystyle z_{b,j} =\sum _{i=1}^{k} \ z_{b_{i,j}}$
\item Let $\displaystyle \Delta z_{j} =z_{a,j} -z_{b,j}$
\end{itemize}

Choose $\displaystyle \phi \left( x_{1} ,x_{2} ,x_{3} ,x_{4}\right) =\frac{x_{2} \Delta z_{1} -x_{1} \Delta z_{2}}{x_{1} -x_{2}} -\frac{x_{4} \Delta z_{3} -x_{3} \Delta z_{4}}{x_{3} -x_{4}}$ and let $\displaystyle C=0$.



Step 5:

We have $\displaystyle \phi \left( g_{1} ,g_{2} ,g_{3} ,g_{4}\right) =est2\left( C_{1} ;C_{2} ;a_{1} ,...,a_{k} ;b_{1} ,...,b_{k}\right) -est2\left( C_{3} ;C_{4} ;a_{1} ,...,a_{k} ;b_{1} ,...,b_{k}\right)$.

Let $\displaystyle C=0$.



\textbf{Claim:} $\displaystyle \phi \left( g_{1} ,g_{2} ,g_{3} ,g_{4}\right)$ is a consistent estimator of $\displaystyle C=0$ if LLM A = LLM B.



Proof:

Assume LLM A = LLM B.

By theorem 6, both $\displaystyle est2\left( C_{1} ;C_{2} ;a_{1} ,...,a_{k} ;b_{1} ,...,b_{k}\right)$ and $\displaystyle est2\left( C_{3} ;C_{4} ;a_{1} ,...,a_{k} ;b_{1} ,...,b_{k}\right)$ converge to $\displaystyle \Delta b$, so $\displaystyle \phi \left( g_{1} ,g_{2} ,g_{3} ,g_{4}\right)$ converges to 0.\qed 



\subsection{Deferred Proofs ($\star$)}



\textbf{Lemma 1:} $\displaystyle g_{j}$ is a consistent estimator of $\displaystyle c_{j}$ for $\displaystyle j=1:m$ as $\displaystyle N^{\left( 1\right)} ,...,N^{\left( w\right)}\rightarrow \infty $.



Proof: As $\displaystyle N^{\left( 1\right)} ,...,N^{\left( w\right)}\rightarrow \infty $, we have $\displaystyle N_{1} ,...,N_{m}\rightarrow \infty $. The lemma then follows from the fact that $\displaystyle \frac{n_{j}\left( a,i\right)}{n_{j}\left( b,i\right)} =\frac{n_{j}\left( a,i\right) /N_{j}}{n_{j}\left( b,i\right) /N_{j}}$ are consistent estimators of $\displaystyle \frac{p_{j}\left( a,i\right)}{p_{j}\left( b,i\right)}$.\qed 



\textbf{Lemma 2:} $\displaystyle \phi \left( g_{1} ,...,g_{m}\right)$ is a consistent estimator of $\displaystyle \phi \left( c_{1} ,...,c_{m}\right)$ as $\displaystyle N^{\left( 1\right)} ,...,N^{\left( w\right)}\rightarrow \infty $.

Proof: Follows from lemma 1.\qed 



\textbf{Lemma 3:} The solution set $\displaystyle S=\left\{x\in \mathbb{R}^{m} :\phi \left( x\right) =C\right\}$ has Lebesgue measure 0.

Proof: $\displaystyle \phi \left( x\right)$ (and therefore $\displaystyle \phi \left( x\right) -C$) is real-analytic by assumption. Furthermore, $\displaystyle \phi \left( x\right) -C$ is not identically zero by assumption. A standard result in real-analytic geometry states that the zero set of a nontrivial real-analytic function on $\displaystyle \mathbb{R}^{m}$ is a countable union of manifolds of dimension at most $\displaystyle m-1$, and therefore has Lebesgue measure zero.\qed 



Now consider $\displaystyle c=\left( c_{1} ,...,c_{m}\right) =\left( ln\left(\prod _{i=1}^{k_{1}}\frac{p_{1}\left( a,i\right)}{p_{1}\left( b,i\right)}\right) ,...,ln\left(\prod _{i=1}^{k_{m}}\frac{p_{m}\left( a,i\right)}{p_{m}\left( b,i\right)}\right)\right)$. Let $\displaystyle \Theta $ denote the configuration of LLM B, including architecture, model weights, and decoding parameters. Let $\displaystyle c\left( \Theta \right)$ (or just $\displaystyle c$ when the context is clear) denote the values of $\displaystyle \left( c_{1} ,...,c_{m}\right)$ corresponding to $\displaystyle \Theta $. When LLM A $\displaystyle \neq $LLM B, we treat $\displaystyle c\left( \Theta \right)$ as a random vector over the distribution of "potential" LLM configurations $\displaystyle \Theta $.



The rest of our proofs require the following axiom:



\textbf{Axiom 1:} If LLM A$\displaystyle \neq $LLM B, then $\displaystyle c\left( \Theta \right) \sim \pi $ where $\displaystyle \pi $ is a probability measure with $\displaystyle \pi \left( S\right) =0$.



Interpretive modeling justification of Axiom 1:

Before we continue, note that nothing below is meant to prove the axiom (as indeed, the axiom is logically independent of our previous assumptions). This is a modeling narrative aiming to justify the axiom as a principled modeling assumption (analogous to the role the Church-Turing Thesis plays in computability theory, which links a formal mathematical condition to an informal notion of real-world behavior and is justified by modeling considerations rather than proof).



Lemma 3 demonstrates that $\displaystyle S$ is a 'pathological' set in the sense of having measure 0, being so small that it takes up 0 measurable space within $\displaystyle \mathbb{R}^{m}$. LLM A = LLM B can be viewed as a pathological special case that collapses $\displaystyle c$ onto this low-dimensional subset $\displaystyle S$.



On the other hand, when LLM A$\displaystyle \neq $LLM B, we are in the 'generic case' where there is no reason, so to speak, for $\displaystyle c$ to collapse onto $\displaystyle S$. In fact, one might go a step further to argue that $\displaystyle c$ should be viewed as an absolutely continuous random vector in $\displaystyle \mathbb{R}^{m}$. Indeed, if assumed, this would imply Axiom 1 as a corollary since any absolutely continuous probability measure with respect to Lebesgue measure necessarily assigns probability 0 to $\displaystyle S$. However, we choose to keep Axiom 1 as our most minimal assumption. It claims that if \ LLM A$\displaystyle \neq $LLM B, then the probability of $\displaystyle c$ landing on the pathological set $\displaystyle S$ is 0.



Implicitly, we've been assuming that when LLM A$\displaystyle \neq $LLM B, LLM B's configuration is drawn from a 'natural' and 'non-adversarial' distribution of potential LLM configurations. This is not the case if the actor behind LLM B artificially manipulates LLM B's configuration in such a precise and specific way to force $\displaystyle c$ to land on $\displaystyle S$. We argue that we should assume that this doesn't happen under our assumed threat model. Why should we assume that the actor behind LLM B doesn't attempt to force $\displaystyle c\in S$ to happen when LLM A$\displaystyle \neq $LLM B? This question concerns the incentives of a rational actor behind LLM B given our threat model. We assert the actor has a disincentive to deliberately configure LLM B in such a way such that $\displaystyle c\in S$ happens.



Suppose LLM A$\displaystyle \neq $LLM B. Note that if $\displaystyle \phi \left( g_{1} ,...,g_{m}\right)$ converges to $\displaystyle C$, both the estimator algorithm and the hypothesis test will erroneously detect LLM B as being LLM A as $\displaystyle N^{\left( 1\right)} ,...,N^{\left( w\right)}\rightarrow \infty $. This isn't obvious for the case of the hypothesis test. The later proof for property 1 of the hypothesis test will show that if $\displaystyle \phi \left( g_{1} ,...,g_{m}\right)$ converges to $\displaystyle C$, then the hypothesis test will decide LLM A = LLM B (i.e. fail to reject $\displaystyle H_{0}$) as $\displaystyle N^{\left( 1\right)} ,...,N^{\left( w\right)}\rightarrow \infty $. Crucially, the proof of property 1 of the hypothesis test will not require assuming Axiom 1 (or else we would have circular reasoning). By lemma 2, $\displaystyle \phi \left( g_{1} ,...,g_{m}\right)$ being a consistent estimator of $\displaystyle C$ is equivalent to $\displaystyle \phi \left( c_{1} ,...,c_{m}\right) =C$. By definition of $\displaystyle S$, $\displaystyle \phi \left( c_{1} ,...,c_{m}\right) =C$ is equivalent to $\displaystyle c\in S$. Therefore, satisfying the condition $\displaystyle c\in S$ causes both the estimator algorithm and hypothesis test to erroneously classify LLM B as being identical to LLM A. Thus, the motive of an actor attempting to configure LLM B in such a way as to force $\displaystyle c\in S$ is to impersonate LLM A with LLM B. Earlier, in the "Setting" section of this paper, we justified why a rational actor would not attempt to impersonate LLM A with LLM B under the threat scenario (and imposed it as an explicit assumption in the threat scenario for good measure). Therefore, we assert that it is justified to assume that the actor behind LLM B does not deliberately attempt to force $\displaystyle c\in S$. 



A final potential objection to this axiom is the possibility of LLM B being a derived version of LLM A, whose implementation difference may be so negligible in principle that $\displaystyle c\in S$ becomes a real possibility with nonzero probability. However, our threat scenario explicitly excluded the possibility of LLM B being a fine-tune or some other derived version of LLM A when LLM A$\displaystyle \neq $LLM B.



Thus, under the threat setting, when LLM B$\displaystyle \neq $LLM A, we assert that LLM B should be viewed as being drawn from a non-adversarial distribution of independently-trained "potential" LLMs. This, combined with the measure 0 size of $\displaystyle S$, concludes our argument for Axiom 1 being a justified modeling assumption for the threat model. \(\circ\)



\textbf{A note about interpreting $\displaystyle \pi$:}

The probability measure $\displaystyle \pi $ in Axiom 1 may be interpreted either in a Bayesian sense (as a prior over potential configurations of LLM B) or in a frequentist sense (as a data-generating distribution over the configurations being created and encountered in practice). Our results are agnostic to which interpretation is used.



Having established Axiom 1, let us now continue with the proofs of the properties of the estimator algorithm and the hypothesis test.



\textbf{Property 1 of the Estimator Algorithm: }$\displaystyle \phi \left( g_{1} ,...,g_{m}\right)$ is a convergent statistic as $\displaystyle N^{\left( 1\right)} ,...,N^{\left( w\right)}\rightarrow \infty $.

Proof: Lemma 2.\qed 



\textbf{Property 2 of the Estimator Algorithm: }If LLM A = LLM B, this $\displaystyle \phi \left( g_{1} ,...,g_{m}\right)$ converges to $\displaystyle C$. Else, it converges to a non-$\displaystyle C$ value with probability 1.

Proof:

The estimator algorithm already assumes that $\displaystyle \phi \left( g_{1} ,...,g_{m}\right)$ converges to $\displaystyle C$ when LLM A = LLM B.



Next, assume LLM A$\displaystyle \neq $LLM B. We wish to show that the probability of $\displaystyle \phi \left( g_{1} ,...,g_{m}\right)$ converging to $\displaystyle C$ is 0. By lemma 2, $\displaystyle \phi \left( g_{1} ,...,g_{m}\right)$ converges to $\displaystyle \phi \left( c_{1} ,...,c_{m}\right)$. Thus, we must show that the probability of $\displaystyle \phi \left( c_{1} ,...,c_{m}\right) =C$ is 0. $\displaystyle \phi \left( c_{1} ,...,c_{m}\right) =C$ is equivalent to $\displaystyle c\in S=\left\{x\in \mathbb{R}^{m} :\phi \left( x\right) =C\right\}$. By axiom 1, the probability that $\displaystyle c\in S$ is 0.\qed 



We now prove the properties of the hypothesis test.



\textbf{Property 1 of the hypothesis test:} If $\displaystyle H_{0}$ if true, with probability 1, the test does not reject $\displaystyle H_{0}$ as $\displaystyle N^{\left( 1\right)} ,...,N^{\left( w\right)}\rightarrow \infty $.

Proof:

Let us assume that $\displaystyle H_{0}$ is true (i.e. that LLM A = LLM B). Let us then establish some facts:

Fact 1: $\displaystyle \phi \left( g_{1} ,...,g_{m}\right)\rightarrow C$ as $\displaystyle N^{\left( 1\right)} ,...,N^{\left( w\right)}\rightarrow \infty $ (by assumption)

Fact 2: $\displaystyle \phi \left( g_{1} ,...,g_{m}\right)$ approaches a normal distribution as $\displaystyle N^{\left( 1\right)} ,...,N^{\left( w\right)}\rightarrow \infty $ (by Theorem 16, which is presented later in the paper)

Fact 3: $\displaystyle \sigma \left( \phi \left( g_{1} ,..,g_{m}\right)\right)$ converges to 0 in $\displaystyle \Theta \left(\sqrt{\sum _{r=1}^{w}\frac{1}{N^{\left( r\right)}}}\right)$ as $\displaystyle N^{\left( 1\right)} ,...,N^{\left( w\right)}\rightarrow \infty $. (a direct consequence of Theorem 17, which is presented later in the paper)



Fact 4: Let $\displaystyle X\sim \mathcal{N}\left( 0,\sigma ^{2}\right)$. Fix a percentile level $\displaystyle p\in \left( 0,1\right)$. Let $\displaystyle c_{p}\left( \sigma \right) =inf\left\{x:P\left( X\leqslant x\right) \geqslant p\right\}$ be the $\displaystyle p$-quantile of $\displaystyle X$. Then $\displaystyle c_{p} =z_{p} \sigma $ where $\displaystyle z_{p} =\Phi ^{-1}\left( p\right)$ and $\displaystyle \Phi $ is the standard normal CDF.



Proof of Fact 4:

$\displaystyle P\left( X\leqslant x\right) =P\left( Z\leqslant \frac{x}{\sigma }\right)$. Thus, $\displaystyle P\left( X\leqslant c_{p}\left( \sigma \right)\right) =p$ (by definition) is equivalent to $\displaystyle P\left( Z\leqslant \frac{c_{p}\left( \sigma \right)}{\sigma }\right) =p$. By definition, we then have $\displaystyle \frac{c_{p}\left( \sigma \right)}{\sigma } =\Phi ^{-1}\left( p\right) \Longrightarrow c_{p}\left( \sigma \right) =z_{p} \sigma $. This completes the proof of fact 4.



Fact 5: Fix a percentile level $\displaystyle p\in \left( 0,1\right)$. Then $\displaystyle c_{p}\left( \sigma \right)$ converges to 0 in $\displaystyle \Theta \left( \sigma \right)$ (direct consequence of fact 4).



From theorem 17 and theorem 19, we have that as $\displaystyle N^{\left( 1\right)} ,...,N^{\left( w\right)}\rightarrow \infty $, the standard deviation of $\displaystyle \phi \left( g_{1} ,...,g_{m}\right)$ dominates its bias (see the comment following theorem 19). We may therefore consider bias to be negligible. As $\displaystyle N^{\left( 1\right)} ,...,N^{\left( w\right)}\rightarrow \infty $, by facts 1-3, we may assume that $\displaystyle \phi \mathcal{\left( g_{1} ,...,g_{m}\right) \sim N}\left( C,\Theta \left(\sqrt{\sum _{r=1}^{w}\frac{1}{N^{\left( r\right)}}}\right)^{2}\right)$.



To prove property 1, it suffices to show that any $\displaystyle p$-quantile of $\displaystyle \phi \left( g_{1} ,...,g_{m}\right)$ converges to $\displaystyle C$ faster than $\displaystyle C\pm \epsilon \left( N^{\left( 1\right)} ,...,N^{\left( w\right)}\right)$ does as $\displaystyle N^{\left( 1\right)} ,...,N^{\left( w\right)}\rightarrow \infty $. Or, equivalently, it suffices to show that any $\displaystyle p$-quantile $\displaystyle c_{p}$ of $\displaystyle \phi \left( g_{1} ,...,g_{m}\right) -C\sim \mathcal{N}\left( 0,\Theta \left(\sqrt{\sum _{r=1}^{w}\frac{1}{N^{\left( r\right)}}}\right)^{2}\right)$ converges to 0 faster than $\displaystyle \pm \epsilon \left( N^{\left( 1\right)} ,...,N^{\left( w\right)}\right)$ does as $\displaystyle N^{\left( 1\right)} ,...,N^{\left( w\right)}\rightarrow \infty $.

Fact 5 tells us that any $\displaystyle p$-quantile $\displaystyle c_{p}$ of $\displaystyle \phi \left( g_{1} ,...,g_{m}\right) -C$ converges to 0 with $\displaystyle \Theta \left(\sqrt{\sum _{r=1}^{w}\frac{1}{N^{\left( r\right)}}}\right)$ speed as $\displaystyle N^{\left( 1\right)} ,...,N^{\left( w\right)}\rightarrow \infty $.

Since our hypothesis test enforces $\displaystyle \frac{\epsilon \left( N^{\left( 1\right)} ,...,N^{\left( w\right)}\right)}{\sqrt{\sum _{r=1}^{w} 1/N^{\left( r\right)}}}\rightarrow \infty $ as $\displaystyle N^{\left( 1\right)} ,...,N^{\left( w\right)}\rightarrow \infty $, indeed, we have that $\displaystyle \pm \epsilon \left( N^{\left( 1\right)} ,...,N^{\left( w\right)}\right)$ converges to 0 slower than $\displaystyle \sqrt{\sum _{r=1}^{w} 1/N^{\left( r\right)}}$ and therefore $\displaystyle c_{p}$ too as $\displaystyle N^{\left( 1\right)} ,...,N^{\left( w\right)}\rightarrow \infty $. This completes the proof of property 1 of the hypothesis test.\qed 





\textbf{Property 2 of the hypothesis test: }If $\displaystyle H_{1}$ is true, with probability 1, the test rejects $\displaystyle H_{0}$ as $\displaystyle N^{\left( 1\right)} ,...,N^{\left( w\right)}\rightarrow \infty $.



Proof:

Assume $\displaystyle H_{1}$ is true (i.e. LLM A$\displaystyle \neq $LLM B). By property 2 of the estimator algorithm, we have that $\displaystyle \phi \left( g_{1} ,...,g_{m}\right)$ converges to a value $\displaystyle C'\neq C$ with probability 1. Let $\displaystyle m=\frac{|C-C'|}{2}$ be the half distance between $\displaystyle C'$ and $\displaystyle C$. Because $\displaystyle \phi \left( g_{1} ,...,g_{m}\right)\rightarrow C'$ as $\displaystyle N^{\left( 1\right)} ,...,N^{\left( w\right)}\rightarrow \infty $, eventually, $\displaystyle \phi \left( g_{1} ,...,g_{m}\right)$ will remain in $\displaystyle \mathbb{R} \backslash \left[ C-m,C+m\right]$. Note that the hypothesis test imposes $\displaystyle \epsilon \left( N^{\left( 1\right)} ,...,N^{\left( w\right)}\right) \downarrow 0$ as $\displaystyle N^{\left( 1\right)} ,...,N^{\left( w\right)}\rightarrow \infty $. Thus, eventually we will have $\displaystyle \left[ C-\epsilon \left( N^{\left( 1\right)} ,...,N^{\left( w\right)}\right) ,C+\epsilon \left( N^{\left( 1\right)} ,...,N^{\left( w\right)}\right)\right] \varsubsetneq \left[ C-m,c+m\right]$ as $\displaystyle N^{\left( 1\right)} ,...,N^{\left( w\right)}\rightarrow \infty $. Since the decision rule is to reject $\displaystyle H_{0}$ if $\displaystyle \phi \left( g_{1} ,...,g_{m}\right) \in \mathbb{R} \backslash \left[ C-m,c+m\right]$, the hypothesis test will eventually reject $\displaystyle H_{0}$ as $\displaystyle N^{\left( 1\right)} ,...,N^{\left( w\right)}\rightarrow \infty $.\qed 



\subsection{Additional remarks for this section}

We've produced \textit{consistent} tests for LLM ownership identity under our setting and assumptions. In this sense, these tests are theoretically "perfect" deciders of whether LLM A = LLM B under our threat scenario. Our general test uses $\displaystyle g_{j}$'s as their fundamental building blocks, and we've explicitly demonstrated that one can create concrete tests from $\displaystyle est$ and $\displaystyle est2$, which themselves are built from these $\displaystyle g_{j}$'s.



\subsubsection{Robustness}

There is another point that we wish to bring attention to. When we built the $\displaystyle \phi \left( g_{1} ,...,g_{m}\right)$'s in the "Specific Examples" subsection, the contexts $\displaystyle C_{j}$ and strings $\displaystyle a_{1,j} ,...,a_{k_{j} ,j} ,b_{1,j} ,...,b_{k_{j} ,j}$ we used to construct them were constrained but otherwise arbitrary choices. This points to the fact that there are a massive class of possible $\displaystyle \phi \left( g_{1} ,...,g_{m}\right)$ estimators that each converge to their corresponding $\displaystyle C$ if LLM A = LLM B, but converge to a non-$\displaystyle C$ value otherwise. The key takeaway is this - there is a vast set of "theoretically perfect" tests for deciding whether or not LLM A = LLM B that can be constructed. The quantity of this redundancy makes testing whether LLM A = LLM B exceptionally robust.



\subsubsection{Generalization to other threat scenarios: Impersonation Attacks}

How well can this test hold up under other threat scenarios? Consider the scenario where LLM A$\displaystyle \neq $LLM B, but the actor behind LLM B is attempting to cause our tests to erroneously detect LLM B as LLM A. In other words, consider the scenario where an impersonation attack were to happen.



For our tests to become asymptotically invalid, it is necessary for Axiom 1 to fail. In other words, such an adversary would need to force $\displaystyle c$ to land on the measure-0 set $\displaystyle S$. Furthermore, if multiple tests are created like in our above robustness discussion, such an adversary would need to force $\displaystyle c$ to collapse onto a measure-0 set $\displaystyle S$ for every single one of these tests. It is worth asking if it is even practically feasible for such an adversary to accomplish this. This motivates our conjecture that it is "hard" for an adversary with a different LLM to impersonate LLM A under our testing scheme, and we encourage further research to explore this direction.



\section{Asymptotic Analysis: Independent $\displaystyle g_{j}$'s}

We begin our analysis of the asymptotic properties of $\displaystyle \phi \left( g_{1} ,...,g_{m}\right)$ by considering a well-behaved but useful case: when $\displaystyle g_{j}$'s are independent. In this section, each $\displaystyle g_{j}$ is assumed to be constructed from an independent, self-contained experiment involving running LLM B many times on context $\displaystyle C_{j}$ and recording number of occurrences of $\displaystyle a_{1,j} ,...,a_{k_{j} ,j} ,b_{1,j} ,...,b_{k_{j} ,j}$.



We will be able to prove precise formulas instead of results written in the more abstract big-$\displaystyle O$ or big-$\displaystyle \Theta $ notation found in the general case (See the next section: "Asymptotic Analysis: Non-independent $\displaystyle g_{j}$'s (The General Case)"). This will allow us to analyze the actual coefficients of the terms, which then allows us to extract higher-resolution insights that are amenable to practical-case analysis, where coefficients often matter. Much of the ideas, tools, techniques, and intuition developed in this section carry forward and help guide the analysis of the generic non-independent case in the final section of this paper.



In this section, we will prove several core results for the independent case, then discuss how these results can be used to design $\displaystyle \phi \left( g_{1} ,...,g_{m}\right)$ estimators that can converge within a 'reasonable' or 'practical' number of calls to LLM B. We will also provide a few examples of variance computations for a couple of concrete estimators of interest.



We will adopt the language and notation of the previous section in this section. When we consider each $\displaystyle g_{j}$ to be constructed from an independent, self-contained experiment, we are saying that the $\displaystyle j$'s and the experiments ($\displaystyle \{E_{r} :r=1:w\}$) have a 1-1 correspondence. Each $\displaystyle j=1:m$ corresponds to a unique $\displaystyle r=1:w$ and vice versa. This implies $\displaystyle w=m$. This also implies that the $\displaystyle N_{j}$'s and $\displaystyle N^{\left( r\right)}$'s are identical up to reordering. It is therefore redundant to speak in the language of experiments in this independent $\displaystyle g_{j}$ setting. For bookkeeping simplicity, we will speak in the language of $\displaystyle j$-indexed concepts rather than in the language of experiments ($\displaystyle r$-indexed concepts) for the rest of this section. Just keep in mind that they are interchangeable concepts in this setting.



Let us now state our assumptions for this section clearly:
\begin{mdframed}
\begin{itemize}
    \item Throughout this section, $\phi(x_1, \dots, x_m)$ will be an arbitrary analytic function unless otherwise stated.
    \item Throughout this section, we shall restrict our attention to the case where the $g_j$'s are decoupled and independent. In other words, each $g_j$ is constructed from an independent, self-contained experiment involving running LLM B many times on context $C_j$ and recording number of occurrences of $a_{1,j}, \dots, a_{k_j,j}, b_{1,j}, \dots, b_{k_j,j}$.
    \item Throughout this section, we will be assuming $\dagger\dagger$ instead of $\dagger$ (which is a special case of $\dagger\dagger$; $\dagger\dagger$ is a relaxation of $\dagger$).
    \item Additional Regularity Condition (*): For each $C_j$, it is impossible for $a_{1,j}, \dots, a_{k_j,j}, b_{1,j}, \dots, b_{k_j,j}$ to be string completions of each other with respect to LLM B (practical case: $a_{1,j}, \dots, a_{k_j,j}, b_{1,j}, \dots, b_{k_j,j}$ are not substrings of each other).
\end{itemize}
\end{mdframed}
Note: (*) imposes a nice regularity structure that will allow us to derive clean, precise variance formulas. In particular, due to (*) preventing simultaneous generations among $\displaystyle a_{1,j} ,...,a_{k_{j} ,j} ,b_{1,j} ,...,b_{k_{j} ,j}$ for each $\displaystyle j$, the $\displaystyle n_{j}\left( a,i\right)$'s and $\displaystyle n_{j}\left( b,i\right)$'s are components of a multinomial distribution for each $\displaystyle j$.



\subsection{The Core Theorems}

As it turns out, $\displaystyle \phi \left( g_{1} ,...,g_{m}\right)$ and the $\displaystyle g_{j}$'s have multiple desirable estimator properties:



\textbf{Theorem 7:} $\displaystyle g_{1} ,...,g_{m}$ are independent.



Proof: By construction.\qed 



\textbf{Theorem 8}: $\displaystyle g_{j}$ is a consistent estimator of $\displaystyle c_{j}$ as $\displaystyle N_{j}\rightarrow \infty $.

Proof: This is a restatement of lemma 1. We've inserted it here for presentation.\qed 



\textbf{Theorem 9:} $\displaystyle var\left( g_{j}\right)\rightarrow \frac{1}{N_{j}}\left(\sum _{i=1}^{k_{j}}\frac{1}{p_{j}\left( a,i\right)} +\sum _{i=1}^{k_{j}}\frac{1}{p_{j}\left( b,i\right)}\right)$ as $\displaystyle N_{j}\rightarrow \infty $

Proof: Deferred for later ($\displaystyle \star \star $).\qed 



Comment: $\displaystyle g_{j}$'s variance converges in $\displaystyle O\left( 1/N_{j}\right)$ speed, which matches the best possible parametric decay rate $\displaystyle O\left( 1/N_{j}\right)$ for the variance of an estimator satisfying standard regularity conditions.



\textbf{Theorem 10:} Each $\displaystyle g_{j}$ is asymptotically a normal random variable as $\displaystyle N_{j}\rightarrow \infty $.



Proof: Deferred for later ($\displaystyle \star \star $).\qed 



\textbf{Theorem 11:} $\displaystyle \phi \left( g_{1} ,...,g_{m}\right)$ is a consistent estimator of $\displaystyle \phi \left( c_{1} ,...,c_{m}\right)$ as $\displaystyle N_{1} ,...,N_{m}\rightarrow \infty $



Proof: Theorem 8.\qed 

\textbf{Theorem 12:} $\displaystyle var\left( \phi \left( g_{1} ,...,g_{m}\right)\right)\rightarrow \sum _{j=1}^{m}\frac{1}{N_{j}}\left(\frac{\partial \phi }{\partial g_{j}}\left( c_{1} ,...,c_{m}\right)\right)^{2}\left(\sum _{i=1}^{k_{j}}\frac{1}{p_{j}\left( a,i\right)} +\sum _{i=1}^{k_{j}}\frac{1}{p_{j}\left( b,i\right)}\right)$ as $\displaystyle N_{1} ,...,N_{m}\rightarrow \infty $.

Proof:

By Theorem 8, $\displaystyle g_{j}$ is a consistent estimator of $\displaystyle c_{j}$, so as $\displaystyle N_{1} ,...,N_{m}\rightarrow \infty $, we have by the delta method:

$\displaystyle \phi \left( g_{1} ,...,g_{m}\right) \approx \phi \left( c_{1} ,...,c_{m}\right) +\sum _{j=1}^{m}\frac{\partial \phi }{\partial g_{j}}\left( c_{1} ,...,c_{m}\right)\left( g_{j} -c_{j}\right)$ 

$\displaystyle \Longrightarrow var\left( \phi \left( g_{1} ,...,g_{m}\right)\right) \approx \sum _{j=1}^{m}\left(\frac{\partial \phi }{\partial g_{j}}\left( c_{1} ,...,c_{m}\right)\right)^{2} var\left( g_{j}\right)$

$\displaystyle =\sum _{j=1}^{m}\frac{1}{N_{j}}\left(\frac{\partial \phi }{\partial g_{j}}\left( c_{1} ,...,c_{m}\right)\right)^{2}\left(\sum _{i=1}^{k_{j}}\frac{1}{p_{j}\left( a,i\right)} +\sum _{i=1}^{k_{j}}\frac{1}{p_{j}\left( b,i\right)}\right)$ (Theorem 9)\qed 



Comment: Since we know $\displaystyle var\left( \phi \left( g_{1} ,...,g_{m}\right)\right)$ is $\displaystyle O\left(\sum _{j=1}^{m}\frac{1}{N_{j}}\right)$ as $\displaystyle N_{1} ,...,N_{m}\rightarrow \infty $, then taking the $\displaystyle N_{j}$'s to be fixed proportions of $\displaystyle N=\sum N_{j}$ (the total number of samples/calls/runs across all experiments) results in $\displaystyle O\left( 1/N\right)$ variance decay for $\displaystyle \phi \left( g_{1} ,...,g_{m}\right)$. This matches the best possible parametric decay rate $\displaystyle O\left( 1/N\right)$ for the variance of an estimator satisfying standard regularity conditions.



\textbf{Theorem 13:} $\displaystyle \phi \left( g_{1} ,...,g_{m}\right)$ is asymptotically a normal random variable as $\displaystyle N_{1} ,...,N_{m}\rightarrow \infty $.



Proof: 

By Theorem 8, $\displaystyle g_{j}$ is a consistent estimator of $\displaystyle c_{j}$, so as $\displaystyle N_{1} ,...,N_{m}\rightarrow \infty $, we have by the delta method:

$\displaystyle \phi \left( g_{1} ,...,g_{m}\right) \approx \phi \left( c_{1} ,...,c_{m}\right) +\sum _{j=1}^{m}\frac{\partial \phi }{\partial g_{j}}\left( c_{1} ,...,c_{m}\right)\left( g_{j} -c_{j}\right)$ 

Theorem 10 tells us that the $\displaystyle g_{j}$'s are asymptotically normal and Theorem 7 tells us that they are independent. Then, the right hand side of the equation is an affine transformation of independent, asymptotically normal random variables. It is a well known fact that an affine transformation of independent normal random variables is normal. This completes the proof. \qed 



\textbf{Theorem 14: }$\displaystyle Bias\left( \phi \left( g_{1} ,...,g_{m}\right)\right) =E\left[ \phi \left( g_{1} ,...,g_{m}\right)\right] -\phi \left( c_{1} ,...,c_{m}\right) =O\left(\sum _{j=1}^{m} 1/N_{j}\right)$

Proof: Deferred for later ($\displaystyle \star \star $).\qed 



Comment: From theorem 12, we have that $\sigma( \phi ( g_{1} ,...,g_{m}))$ decays at rate $\Theta (\sqrt{\sum _{j=1}^{m} 1/N_{j}})$, which asymptotically dominates $Bias( \phi ( g_{1} ,...,g_{m})) =O(\sum _{j=1}^{m} 1/N_{j})$. In other words, $Bias( \phi ( g_{1} ,...,g_{m})) =o( \sigma ( \phi ( g_{1} ,...,g_{m}))$. Combined with the fact that $\displaystyle \phi ( g_{1} ,...,g_{m})$ is asymptotically normal, this allows us to treat bias as negligible ($\displaystyle =0$) as $\displaystyle N_{1} ,...,N_{m}\rightarrow \infty $ for the rest of this section.



\subsubsection*{Heuristic Analysis}

As we've discussed before, $\displaystyle \phi \left( g_{1} ,...,g_{m}\right)$ is an estimator that serves as the central tool in our tests for LLM ownership identity. Naturally, we desire to optimize its convergence speed. Since $\displaystyle \phi \left( g_{1} ,...,g_{m}\right)$ is consistent and asymptotically normal, this problem reduces to minimizing its asymptotic variance.



Let us inspect the equation $\displaystyle var\left( \phi \left( g_{1} ,...,g_{m}\right)\right)\rightarrow \sum _{j=1}^{m}\frac{1}{N_{j}}\left(\frac{\partial \phi }{\partial g_{j}}\left( c_{1} ,...,c_{m}\right)\right)^{2}\left(\sum _{i=1}^{k_{j}}\frac{1}{p_{j}\left( a,i\right)} +\sum _{i=1}^{k_{j}}\frac{1}{p_{j}\left( b,i\right)}\right)$. It is crucial to note that we do not, in practice, know the values of the $\displaystyle p_{j}\left( a,i\right)$'s and $\displaystyle p_{j}\left( b,i\right)$'s. After all, they are based on black-box LLM B's internal configuration, which we do not assume access to. However, this does not mean we cannot extract heuristic insights from the structure of this equation.



First, note that heuristically speaking, $\displaystyle var\left( \phi \left( g_{1} ,...,g_{m}\right)\right)$ seems to increase as $\displaystyle m$ and the $\displaystyle k_{j}$'s increase (assuming all else being held equal, particularly the values of $\displaystyle \frac{\partial \phi }{\partial g_{j}}\left( c_{1} ,...,c_{m}\right)$). It therefore seems reasonable to try to minimize $\displaystyle m$ and the $\displaystyle k_{j}$'s when building the estimator $\displaystyle \phi \left( g_{1} ,...,g_{m}\right)$, as well as maximize the $\displaystyle p_{j}\left( a,i\right)$ and $\displaystyle p_{j}\left( b,i\right)$'s.



Note that $\displaystyle est$ and $\displaystyle est2$ themselves are each built from 2 instances of $\displaystyle g_{j}$. Therefore, any estimator $\displaystyle \phi \left( g_{1} ,...,g_{m}\right)$ built using $\displaystyle est$ or $\displaystyle est2$ as one of its components has $\displaystyle m\geqslant 2$. Naive Method I and Naive Method II each has $\displaystyle m=4$.



For maximizing the $\displaystyle p_{j}\left( a,i\right)$ and $\displaystyle p_{j}\left( b,i\right)$'s, again, we do not have access to LLM B's internals, so we cannot do this for certain. However, we can use our human semantic intuition about LLM outputs to intelligently guess high-probability strings $\displaystyle a_{i,j} ,b_{i,j}$ for LLM B on context $\displaystyle C_{j}$. For example, if $\displaystyle C_{j}$ were:

Then we might choose $\displaystyle a_{i,j}$ to be "0" and $\displaystyle b_{i,j}$ to be "1".



We summarize our heuristic takeaways with 3 rules of thumb for optimizing the convergence speed of $\displaystyle \phi \left( g_{1} ,...,g_{m}\right)$:
\begin{enumerate}
\item Try to make $\displaystyle m$ small.
\item Try to make the $\displaystyle k_{j}$'s small (aim for $\displaystyle k_{j} =1$ for each $\displaystyle j=1:m$).
\item Try to make each $\displaystyle p_{j}\left( a,i\right)$ and $\displaystyle p_{j}\left( b,i\right)$ large by intelligently choosing strings $\displaystyle a_{i,j}$ and $\displaystyle b_{i,j}$ for context $\displaystyle C_{j}$.
\end{enumerate}



\subsection{Optimizing convergence speed under a fixed number-of-runs budget}

Consider again the equation $\displaystyle var\left( \phi \left( g_{1} ,...,g_{m}\right)\right)\rightarrow \sum _{j=1}^{m}\frac{1}{N_{j}}\left(\frac{\partial \phi }{\partial g_{j}}\left( c_{1} ,...,c_{m}\right)\right)^{2}\left(\sum _{i=1}^{k_{j}}\frac{1}{p_{j}\left( a,i\right)} +\sum _{i=1}^{k_{j}}\frac{1}{p_{j}\left( b,i\right)}\right)$. One thing we haven't addressed yet is how to choose the $\displaystyle N_{j}$'s to maximize convergence speed. In practice, our number-of-runs budget is constrained as $\displaystyle N_{1} +...+N_{m} =N$ for some fixed $\displaystyle N$. A method for computing optimal values of $\displaystyle N_{1} ,...,N_{m}$ given a compute budget of $\displaystyle N$ runs is given below, which assumes that we have the values of the $\displaystyle p_{j}\left( a,i\right)$'s and $\displaystyle p_{j}\left( b,i\right)$'s (Again, in practice, this is a nontrivial assumption. We will discuss a heuristic method for computing optimal values of $\displaystyle N_{j}$ in the "Heuristic Proportion Selection" subsection).



Let us introduce a new shorthand:
\begin{itemize}
\item Let $\displaystyle q_{j} =\left(\frac{\partial \phi }{\partial g_{j}}\left( c_{1} ,...,c_{m}\right)\right)^{2}\left(\sum _{i=1}^{k_{j}}\frac{1}{p_{j}\left( a,i\right)} +\sum _{i=1}^{k_{j}}\frac{1}{p_{j}\left( b,i\right)}\right)$. Note that $\displaystyle q_{j}$ is a nonnegative constant.
\end{itemize}



\textbf{Theorem 15:} Assume $\displaystyle N_{1} ,...,N_{m}\rightarrow \infty $. Given a run budget of $\displaystyle N_{1} +...+N_{m} =N$, $\displaystyle N_{j}^{*} \approx \frac{N\sqrt{q_{j}}}{\sum _{l=1}^{m}\sqrt{q_{l}}}$ becomes the optimal choice of $\displaystyle N_{j}$ to maximize convergence speed, with $\displaystyle var\left( \phi \left( g_{1} ,...,g_{m}\right)\right)\rightarrow \frac{\left(\sum _{j=1}^{m}\sqrt{q_{j}}\right)^{2}}{N}$ in this case.



Proof:

By theorem 12, we have that $\displaystyle var\left( \phi \left( g_{1} ,...,g_{m}\right)\right)\rightarrow \sum _{j=1}^{m}\frac{q_{j}}{N_{j}}$. We apply Lagrange multipliers to the problem of optimizing $\displaystyle \sum _{j=1}^{m}\frac{q_{j}}{N_{j}}$ under the constraint $\displaystyle N_{1} +...+N_{m} =N$. Note that we intentionally did not include the $\displaystyle N_{1} ,...,N_{m} \geqslant 0$ constraints. We will later find that the optimizer actually satisfies those constraints anyway.



Lagrangian: $\displaystyle \mathcal{L}\left( N_{1} ,...,N_{m} ;\lambda \right) =\sum _{j=1}^{m}\frac{q_{j}}{N_{j}} +\lambda \left(\sum _{j=1}^{m} N_{j} -N\right)$.

$\displaystyle \frac{\partial \mathcal{L}}{\partial N_{j}} =0\Longrightarrow \lambda =\frac{q_{j}}{N_{j}^{2}} \Longrightarrow N_{j} =\sqrt{\frac{q_{j}}{\lambda }}$ (2)

By the constraint $\displaystyle \sum N_{j} =N$, we have:

$\displaystyle \sum _{j=1}^{m}\sqrt{\frac{q_{j}}{\lambda }} =N$

$\displaystyle \Longrightarrow \lambda =\frac{\left(\sum _{j=1}^{m}\sqrt{q_{j}}\right)^{2}}{N^{2}}$ (3)

Plugging (3) into (2) to get $\displaystyle N_{j}^{*} =\frac{N\sqrt{q_{j}}}{\sum _{l=1}^{m}\sqrt{q_{l}}}$.

Observe that $\displaystyle N_{j}^{*}  >0$, which satisfies the implicit $\displaystyle N_{1} ,...,N_{m} \geqslant 0$ constraints. Furthermore, $\displaystyle N_{j}^{*}\rightarrow \infty $ when $\displaystyle N_{1} +...+N_{m} =N\rightarrow \infty $. Therefore, this optimizer is consistent with our $\displaystyle N_{1} ,...,N_{m}\rightarrow \infty $ assumption. We conclude that $\displaystyle N_{j}^{*} \approx \frac{N\sqrt{q_{j}}}{\sum _{l=1}^{m}\sqrt{q_{l}}}$ is our desired optimizer.



Plugging the $\displaystyle N_{j}^{*}$'s back in, we get:

$\displaystyle  \begin{array}{l}
var\left( \phi \left( g_{1} ,...,g_{m}\right)\right)\rightarrow \sum _{j=1}^{m}\frac{q_{j}}{N_{j}^{*}}\\
=\sum _{j=1}^{m}\frac{q_{j}}{\frac{N\sqrt{q_{j}}}{\sum _{l=1}^{m}\sqrt{q_{l}}}}\\
=\frac{\sum _{l=1}^{m}\sqrt{q_{l}}}{N}\sum _{j=1}^{m} q_{j}\\
=\frac{\left(\sum _{j=1}^{m}\sqrt{q_{j}}\right)^{2}}{N}
\end{array}$

\qed 



Comment: A key structural insight from this theorem is that the optimal values $\displaystyle N_{j}^{*} ,j=1:m$ are fixed proportions of $\displaystyle N$ which don't depend on $\displaystyle N$. Thus, it often makes more practical sense to solve for the proportions rather than solve for $\displaystyle N_{j}^{*}$'s themselves; solving for proportions will allow us to reuse them across different $\displaystyle N$'s.



\textbf{Corollary 5}: Assume $\displaystyle N_{1} ,...,N_{m}\rightarrow \infty $. Given a run budget of $\displaystyle N_{1} +...+N_{m} =N$, we have:
\begin{enumerate}
\item $\displaystyle \frac{\partial }{\partial N_{j}} var\left( \phi \left( g_{1} ,...,g_{m}\right)\right)\rightarrow -\frac{\left(\sum _{j=1}^{m}\sqrt{q_{j}}\right)^{2}}{N^{2}}$ (where we assume $\displaystyle N_{j} =\frac{N\sqrt{q_{j}}}{\sum _{l=1}^{m}\sqrt{q_{l}}}$ for $\displaystyle j=1:m$)
\item Let $\displaystyle F\left( N\right) =min_{\sum N_{j} =N} var\left( \phi \left( g_{1} ,...,g_{m}\right)\right)$. Then $\displaystyle \frac{\partial }{\partial N} F\left( N\right)\rightarrow -\frac{\left(\sum _{j=1}^{m}\sqrt{q_{j}}\right)^{2}}{N^{2}}$
\end{enumerate}



Proof:

Revisiting the Lagrangian of the problem of optimizing $\displaystyle \sum _{j=1}^{m}\frac{q_{j}}{N_{j}}$ under the constraint $\displaystyle N_{1} +...+N_{m} =N$, we have: $\displaystyle \mathcal{L}\left( N_{1} ,...,N_{m} ;\lambda \right) =\sum _{j=1}^{m}\frac{q_{j}}{N_{j}} +\lambda \left(\sum _{j=1}^{m} N_{j} -N\right)$.

$\displaystyle \frac{\partial \mathcal{L}}{\partial N_{j}} =0\Longrightarrow \frac{\partial }{\partial N_{j}} var\left( \phi \left( g_{1} ,...,g_{m}\right)\right) \approx \frac{\partial }{\partial N_{j}}\left(\sum _{j=1}^{m}\frac{q_{j}}{N_{j}}\right) =-\lambda $

Plugging in $\displaystyle \lambda =\frac{\left(\sum _{j=1}^{m}\sqrt{q_{j}}\right)^{2}}{N^{2}}$ from equation (3) proves 1.



2. is a direct consequence of the fact that $\displaystyle var\left( \phi \left( g_{1} ,...,g_{m}\right)\right)$ asymptotically converges to $\displaystyle \sum _{j=1}^{m}\frac{q_{j}}{N_{j}}$ and the Envelope Theorem, a well-known result in constrained optimization\qed 



\subsubsection{Heuristic Proportion Selection}

First, define $\displaystyle s_{j} :=\frac{\sqrt{q_{j}}}{\sum _{l=1}^{m}\sqrt{q_{l}}}$ for $\displaystyle j=1:m$. The core practical issue we continue to face is that the value of $\displaystyle s_{j}$ depends on our knowledge of the $\displaystyle p_{j}\left( a,i\right)$'s and $\displaystyle p_{j}\left( b,i\right)$'s, which requires knowledge of black-box LLM B's internal configuration (not assumed). Below, we outline a heuristic approach for estimating good values of $\displaystyle s_{j}$ for optimizing the convergence speed of $\displaystyle \phi \left( g_{1} ,...,g_{m}\right)$.



\textbf{Definition:} Let $\displaystyle p$ be a vector whose entries are all of the $\displaystyle p_{j}\left( a,i\right)$'s and $\displaystyle p_{j}\left( b,i\right)$'s. In particular, $\displaystyle p$ is a vector in $\displaystyle \mathbb{R}^{\sum 2k_{j}}$.



\textbf{A Heuristic Monte Carlo method:}

\textbf{Goal: }Find a convergence-speed-optimal proportion vector $\displaystyle s=\langle s_{1} ,...,s_{m} \rangle $ representing the proportions of $\displaystyle N$ used to compute the $\displaystyle N_{j}^{*}$'s.

\textbf{Steps:}
\begin{enumerate}
\item Viewing the components of $\displaystyle p$ as random variables, guess a prior distribution $\displaystyle \pi $ for $\displaystyle p$ (e.g. guess a plausible range for each $\displaystyle p_{j}\left( a,i\right)$ and each $\displaystyle p_{j}\left( b,i\right)$ to independently uniformly randomly sample from)
\item Repeatedly sample $\displaystyle p$ from $\displaystyle \pi $ and compute $\displaystyle s=\langle s_{1} ,..,s_{m} \rangle $ for each sample. Note that the computations of the $\displaystyle s$'s are independent of $\displaystyle N$.
\item Take the mean of the $\displaystyle s$'s to get an approximation for $\displaystyle E_{p\sim \pi }\left[ s\right]$. The final vector will be the proportions of $\displaystyle N$ that we can use to compute $\displaystyle N_{j}^{*}$'s.
\end{enumerate}



\textbf{A more sophisticated method for creating a generative }$\displaystyle \pi $\textbf{ distribution:}
\begin{enumerate}
\item Procure a list of white-box LLMs (which may include LLM A and open-source/open-weight models). Call these LLMs $\displaystyle M_{1} ,...,M_{l}$.
\item For each standard token-distribution-modifying decoding parameter, guess a pragmatic range that it lives in.
\item A LLM configuration $\displaystyle \Theta $ is defined as a choice of LLM and exact values for its standard token-distribution-modifying decoding parameters. Sample $\displaystyle \Theta $ by uniformly randomly selecting an LLM from $\displaystyle \left\{M_{1} ,...,M_{l}\right\}$ and standard token-distribution-modifying decoding parameters from their corresponding ranges.
\item For each $\displaystyle \Theta $ sampled, compute the $\displaystyle p_{j}\left( a,i\right)$'s and $\displaystyle p_{j}\left( b,i\right)$'s associated with it.
\end{enumerate}

The overall idea is that although we don't know LLM B's exact $\displaystyle p_{j}\left( a,i\right)$'s and $\displaystyle p_{j}\left( b,i\right)$'s, we can use \textit{other} LLM configurations to compute noisy approximations for them.



\subsection{Bringing it all together: A Full Pipeline}

Let us aggregate our previous ideas into a full pipeline for $\displaystyle \phi \left( g_{1} ,...,g_{m}\right)$ estimator design and estimation.



Note: There is a classic optimization/pedagogical clarity tradeoff when it comes to presenting algorithms. We have chosen to present this pipeline in an illustrative, conceptually clear way. However, note that for a practical deployment, there will be redundant computations in the pseudocode below if it is implemented exactly as-is.



\textbf{Goal:} Produce an estimator $\displaystyle \phi \left( g_{1} ,...,g_{m}\right)$ that converges quickly.



\textbf{Steps:}
\begin{enumerate}
\item Using our 3 heuristic rules of thumb, construct a list of candidate estimators $\displaystyle \phi \left( g_{1} ,...,g_{m}\right)$. Do not yet instantiate the values of the $\displaystyle N_{j}$'s.
\item For each candidate estimator $\displaystyle \phi \left( g_{1} ,...,g_{m}\right)$, compute a proportion vector $\displaystyle s_{\phi \left( g_{1} ,...,g_{m}\right)} =\langle s_{1} ,..,s_{m} \rangle $ using heuristic proportion selection.
\item For each candidate estimator $\displaystyle \phi \left( g_{1} ,...,g_{m}\right)$ and its corresponding proportion vector proportion vector $\displaystyle s_{\phi \left( g_{1} ,...,g_{m}\right)} =\langle s_{1} ,..,s_{m} \rangle $, compute $\displaystyle v_{\phi \left( g_{1} ,...,g_{m}\right)} =\sum _{j=1}^{m}\frac{1}{s_{j}}\left(\frac{\partial \phi }{\partial g_{j}}\left( c_{1} ,...,c_{m}\right)\right)^{2}\left(\sum _{i=1}^{k_{j}}\frac{1}{p_{j}\left( a,i\right)} +\sum _{i=1}^{k_{j}}\frac{1}{p_{j}\left( b,i\right)}\right)$.
\item Choose the candidate estimator $\displaystyle \phi \left( g_{1} ,...,g_{m}\right)$ and corresponding proportion vector $\displaystyle s_{\phi \left( g_{1} ,...,g_{m}\right)} =\langle s_{1} ,..,s_{m} \rangle $ with the lowest $\displaystyle v_{\phi \left( g_{1} ,...,g_{m}\right)}$.
\end{enumerate}



Note: $\displaystyle s_{\phi \left( g_{1} ,...,g_{m}\right)}$ is a heuristically chosen proportion vector, representing our best guess of the proportions of $\displaystyle N$ that each $\displaystyle N_{j}^{*}$ must take to minimize variance/convergence speed of $\displaystyle \phi \left( g_{1} ,...,g_{m}\right)$. Based on theorem 12 and theorem 15, $\displaystyle v_{\phi \left( g_{1} ,...,g_{m}\right)}$ is the variance approximation of $\displaystyle \phi \left( g_{1} ,...,g_{m}\right)$ using $\displaystyle s_{\phi \left( g_{1} ,...,g_{m}\right)}$, up to the factor $\displaystyle 1/N$.



Having chosen $\displaystyle \phi \left( g_{1} ,...,g_{m}\right)$ and $\displaystyle s_{\phi \left( g_{1} ,...,g_{m}\right)} =\langle s_{1} ,..,s_{m} \rangle $, we now have a natural estimation algorithm:



\textbf{Goal:} Estimate the limit of $\displaystyle \phi \left( g_{1} ,...,g_{m}\right)$.

\textbf{Steps:}
\begin{enumerate}
\item Iteratively increase the number of total observations $\displaystyle N$ while maintaining $\displaystyle N_{1} ,...,N_{m}$ to have the fixed relative proportions $\displaystyle s_{\phi \left( g_{1} ,...,g_{m}\right)} =\langle s_{1} ,..,s_{m} \rangle $.
\item Apply an early stopping rule.
\end{enumerate}



\subsection{Variance Calculation Examples}

We'll now explore specific variance examples that may of be interest to the reader and which will serve to elucidate some of the previous results.



\subsubsection{The estimator est}
\begin{mdframed}
For the rest of this "The estimator est" subsection, we are assuming that contexts $C_1, C_2$ and strings $a_1, \dots, a_k$ and $b_1, \dots, b_k$ satisfy the ($\dagger$) regularity conditions, in addition to the prevailing assumptions being made throughout the entire "Asymptotic Analysis: Independent $g_j$'s" section.
\end{mdframed}

Recall $\displaystyle est$: $\displaystyle est\left( C_{1} ;C_{2} ;a_{1} ,...,a_{k} ;b_{1} ,...,b_{k}\right) =\frac{\Delta z_{1} -\Delta z_{2}}{ln\left(\prod _{i=1}^{k}\frac{n_{1}\left( a,i\right)}{n_{1}\left( b,i\right)}\right) -ln\left(\prod _{i=1}^{k}\frac{n_{2}\left( a,i\right)}{n_{2}\left( b,i\right)}\right)}$



We can rewrite this as $\displaystyle est\left( C_{1} ;C_{2} ;a_{1} ,...,a_{k} ;b_{1} ,...,b_{k}\right) =\phi \left( g_{1} ,g_{2}\right)$ where $\displaystyle g_{1} =ln\left(\prod _{i=1}^{k}\frac{n_{1}\left( a,i\right)}{n_{1}\left( b,i\right)}\right) ,g_{2} =ln\left(\prod _{i=1}^{k}\frac{n_{2}\left( a,i\right)}{n_{2}\left( b,i\right)}\right) ,\phi \left( g_{1} ,g_{2}\right) =\frac{\Delta z_{1} -\Delta z_{2}}{g_{1} -g_{2}}$.



Note: $\displaystyle est\left( C_{1} ;C_{2} ;a_{1} ,...,a_{k} ;b_{1} ,...,b_{k}\right)$ converges to $\displaystyle T$ when LLM A = LLM B, which is an unknown constant. Therefore, $\displaystyle est\left( C_{1} ;C_{2} ;a_{1} ,...,a_{k} ;b_{1} ,...,b_{k}\right)$ cannot be used as the estimator in the estimator algorithm. Nevertheless, it can be expressed in the form of $\displaystyle \phi \left( g_{1} ,g_{2}\right)$ and the theorems in the previous section can still be used to analyze it.



We then have that theorem 11, theorem 12, and theorem 13 hold, so we know that $\displaystyle est\left( C_{1} ;C_{2} ;a_{1} ,...,a_{k} ;b_{1} ,...,b_{k}\right)$ has desirable properties like consistency, asymptotic normality, and an asymptotic approximation for its variance. Let's compute its asymptotic variance.



By Theorem 12, we have that $\displaystyle var\left( est\left( C_{1} ;C_{2} ;a_{1} ,...,a_{k} ;b_{1} ,...,b_{k}\right)\right)\rightarrow \sum _{j=1}^{2}\frac{1}{N_{j}}\left(\frac{\partial \phi }{\partial g_{j}}\left( c_{1} ,c_{2}\right)\right)^{2}\left(\sum _{i=1}^{k}\frac{1}{p_{j}\left( a,i\right)} +\sum _{i=1}^{k}\frac{1}{p_{j}\left( b,i\right)}\right)$ as $\displaystyle N_{1} ,N_{2}\rightarrow \infty $.

We have:

$\displaystyle \frac{\partial \phi }{\partial g_{1}}\left( c_{1} ,c_{2}\right) =-\frac{\Delta z_{1} -\Delta z_{2}}{\left( c_{1} -c_{2}\right)^{2}}$ and $\displaystyle \frac{\partial \phi }{\partial g_{2}}\left( c_{1} ,c_{2}\right) =\frac{\Delta z_{1} -\Delta z_{2}}{\left( c_{1} -c_{2}\right)^{2}}$.



Thus, we have:

$\displaystyle  \begin{array}{l}
\sum _{j=1}^{2}\frac{1}{N_{j}}\left(\frac{\partial \phi }{\partial g_{j}}\left( c_{1} ,c_{2}\right)\right)^{2}\left(\sum _{i=1}^{k}\frac{1}{p_{j}\left( a,i\right)} +\sum _{i=1}^{k}\frac{1}{p_{j}\left( b,i\right)}\right)\\
=\frac{\left( \Delta z_{1} -\Delta z_{2}\right)^{2}}{\left( c_{1} -c_{2}\right)^{4}}\left(\frac{1}{N_{1}}\left(\sum _{i=1}^{k}\frac{1}{p_{1}\left( a,i\right)} +\sum _{i=1}^{k}\frac{1}{p_{1}\left( b,i\right)}\right) +\frac{1}{N_{2}}\left(\sum _{i=1}^{k}\frac{1}{p_{2}\left( a,i\right)} +\sum _{i=1}^{k}\frac{1}{p_{2}\left( b,i\right)}\right)\right)
\end{array}$



This gives us:



\textbf{Corollary 6:}

$\displaystyle  \begin{array}{l}
var\left( est\left( C_{1} ;C_{2} ;a_{1} ,...,a_{k} ;b_{1} ,...,b_{k}\right)\right)\\
\rightarrow \frac{\left( \Delta z_{1} -\Delta z_{2}\right)^{2}}{\left( c_{1} -c_{2}\right)^{4}}\left(\frac{1}{N_{1}}\left(\sum _{i=1}^{k}\frac{1}{p_{1}\left( a,i\right)} +\sum _{i=1}^{k}\frac{1}{p_{1}\left( b,i\right)}\right) +\frac{1}{N_{2}}\left(\sum _{i=1}^{k}\frac{1}{p_{2}\left( a,i\right)} +\sum _{i=1}^{k}\frac{1}{p_{2}\left( b,i\right)}\right)\right)
\end{array}$ as $\displaystyle N_{1} ,N_{2}\rightarrow \infty $\qed 



Recall that when LLM A = LLM B, we have that $\displaystyle c_{j} =\frac{\Delta z_{j} +\Delta b}{T}$. This gives us:



\textbf{Corollary 7:} If LLM A = LLM B, $var\left( est\left( C_{1} ;C_{2} ;a_{1} ,...,a_{k} ;b_{1} ,...,b_{k}\right)\right)$\\
$\rightarrow \frac{T^{4}}{\left( \Delta z_{1} -\Delta z_{2}\right)^{2}}\left(\frac{1}{N_{1}}\left(\sum _{i=1}^{k}\frac{1}{p_{1}\left( a,i\right)} +\sum _{i=1}^{k}\frac{1}{p_{1}\left( b,i\right)}\right) +\frac{1}{N_{2}}\left(\sum _{i=1}^{k}\frac{1}{p_{2}\left( a,i\right)} +\sum _{i=1}^{k}\frac{1}{p_{2}\left( b,i\right)}\right)\right)$



Proof:

$\displaystyle \frac{\left( \Delta z_{1} -\Delta z_{2}\right)^{2}}{\left( c_{1} -c_{2}\right)^{4}} =\frac{\left( \Delta z_{1} -\Delta z_{2}\right)^{2}}{\left(\frac{\Delta z_{1} +\Delta b}{T} -\frac{\Delta z_{2} +\Delta b}{T}\right)^{4}} =\frac{T^{4}}{\left( \Delta z_{1} -\Delta z_{2}\right)^{2}}$ \qed 



\subsubsection*{Heuristic Analysis:}

Looking at the case of LLM A = LLM B gives us some further heuristic intuition for boosting the convergence speed of $\displaystyle est$. In particular, focus on the $\displaystyle \left( \Delta z_{1} -\Delta z_{2}\right)^{2}$ term. Variance is inversely proportional to it. Recalling our definitions of $\displaystyle \Delta z_{1} ,\Delta z_{2}$, we see that $\displaystyle |\Delta z_{1} -\Delta z_{2} |=|\left(\sum _{i=1}^{k} \ z_{a_{i} ,1} -\sum _{i=1}^{k} \ z_{b_{i} ,1}\right) -\left(\sum _{i=1}^{k} \ z_{a_{i} ,2} -\sum _{i=1}^{k} \ z_{b_{i} ,2}\right) |$.



Fortunately, because we have white-box access to LLM A and therefore all $\displaystyle z_{a_{i} ,j}$'s and $\displaystyle z_{b_{i} ,j}$'s are knowable to us, we can cycle through various candidates for $\displaystyle C_{1} ,C_{2}$ and $\displaystyle a_{1} ,...,a_{k} ,b_{1} ,...,b_{k}$ until our computed value $\displaystyle |\Delta z_{1} -\Delta z_{2} |$ is reasonably large.



To reduce time in producing candidates for $\displaystyle C_{1} ,C_{2}$, and $\displaystyle a_{1} ,...,a_{k} ,b_{1} ,...,b_{k}$, we may adopt the following heuristic: "Pick $\displaystyle C_{1} ,C_{2}$ and $\displaystyle a_{1} ,...,a_{k} ,b_{1} ,...,b_{k}$ such that $\displaystyle a_{1} ,...,a_{k}$ are more likely to be generated than $\displaystyle b_{1} ,...,b_{k}$ on context $\displaystyle C_{1}$ and $\displaystyle b_{1} ,...,b_{k}$ are more likely to be generated than $\displaystyle a_{1} ,...,a_{k}$ on context $\displaystyle C_{2}$ (or vice versa)."



\subsubsection{The estimator est2}
\begin{mdframed}
For the rest of this "The estimator est2" subsection, we are assuming that contexts $C_1, C_2$ and strings $a_1, \dots, a_k$ and $b_1, \dots, b_k$ satisfy the ($\dagger$) regularity conditions, in addition to the prevailing assumptions being made throughout the entire "Asymptotic Analysis: Independent $g_j$'s" section.
\end{mdframed}
Recall $\displaystyle est2$: $\displaystyle est2\left( C_{1} ;C_{2} ;a_{1} ,...,a_{k} ;b_{1} ,...,b_{k}\right) =\frac{ln\left(\prod _{i=1}^{k}\frac{n_{2}\left( a,i\right)}{n_{2}\left( b,i\right)}\right) \Delta z_{1} -ln\left(\prod _{i=1}^{k}\frac{n_{1}\left( a,i\right)}{n_{1}\left( b,i\right)}\right) \Delta z_{2}}{ln\left(\prod _{i=1}^{k}\frac{n_{1}\left( a,i\right)}{n_{1}\left( b,i\right)}\right) -ln\left(\prod _{i=1}^{k}\frac{n_{2}\left( a,i\right)}{n_{2}\left( b,i\right)}\right)}$



We can rewrite this as $\displaystyle est2\left( C_{1} ;C_{2} ;a_{1} ,...,a_{k} ;b_{1} ,...,b_{k}\right) =\phi \left( g_{1} ,g_{2}\right)$ where $\displaystyle g_{1} =ln\left(\prod _{i=1}^{k}\frac{n_{1}\left( a,i\right)}{n_{1}\left( b,i\right)}\right) ,g_{2} =ln\left(\prod _{i=1}^{k}\frac{n_{2}\left( a,i\right)}{n_{2}\left( b,i\right)}\right) ,\phi \left( g_{1} ,g_{2}\right) =\frac{g_{2} \Delta z_{1} -g_{1} \Delta z_{2}}{g_{1} -g_{2}}$.



Note: $\displaystyle est2\left( C_{1} ;C_{2} ;a_{1} ,...,a_{k} ;b_{1} ,...,b_{k}\right)$ converges to $\displaystyle \Delta b$ when LLM A = LLM B, which is an unknown constant. Therefore, $\displaystyle est2\left( C_{1} ;C_{2} ;a_{1} ,...,a_{k} ;b_{1} ,...,b_{k}\right)$ cannot be used as the estimator in the estimator algorithm. Nevertheless, it can be expressed in the form of $\displaystyle \phi \left( g_{1} ,g_{2}\right)$ and the theorems in the previous section can still be used to analyze it.



As usual, we then have that theorem 11, theorem 12, and theorem 13 hold, so we know that $\displaystyle est2\left( C_{1} ;C_{2} ;a_{1} ,...,a_{k} ;b_{1} ,...,b_{k}\right)$ has desirable properties like consistency, asymptotic normality, and an asymptotic approximation for its variance. Let's compute its asymptotic variance.



By Theorem 12, we have that $\displaystyle var\left( est2\left( C_{1} ;C_{2} ;a_{1} ,...,a_{k} ;b_{1} ,...,b_{k}\right)\right)\rightarrow \sum _{j=1}^{2}\frac{1}{N_{j}}\left(\frac{\partial \phi }{\partial g_{j}}\left( c_{1} ,c_{2}\right)\right)^{2}\left(\sum _{i=1}^{k}\frac{1}{p_{j}\left( a,i\right)} +\sum _{i=1}^{k}\frac{1}{p_{j}\left( b,i\right)}\right)$ as $\displaystyle N_{1} ,N_{2}\rightarrow \infty $.

We have:

$\displaystyle \frac{\partial \phi }{\partial g_{1}}\left( c_{1} ,c_{2}\right) =-\frac{c_{2}\left( \Delta z_{1} -\Delta z_{2}\right)}{\left( c_{1} -c_{2}\right)^{2}}$ and $\displaystyle \frac{\partial \phi }{\partial g_{2}}\left( c_{1} ,c_{2}\right) =\frac{c_{1}\left( \Delta z_{1} -\Delta z_{2}\right)}{\left( c_{1} -c_{2}\right)^{2}}$.



Thus, we have:

$\displaystyle  \begin{array}{l}
\sum _{j=1}^{2}\frac{1}{N_{j}}\left(\frac{\partial \phi }{\partial g_{j}}\left( c_{1} ,c_{2}\right)\right)^{2}\left(\sum _{i=1}^{k}\frac{1}{p_{j}\left( a,i\right)} +\sum _{i=1}^{k}\frac{1}{p_{j}\left( b,i\right)}\right)\\
=\frac{\left( \Delta z_{1} -\Delta z_{2}\right)^{2}}{\left( c_{1} -c_{2}\right)^{4}}\left(\frac{c_{2}^{2}}{N_{1}}\left(\sum _{i=1}^{k}\frac{1}{p_{1}\left( a,i\right)} +\sum _{i=1}^{k}\frac{1}{p_{1}\left( b,i\right)}\right) +\frac{c_{1}^{2}}{N_{2}}\left(\sum _{i=1}^{k}\frac{1}{p_{2}\left( a,i\right)} +\sum _{i=1}^{k}\frac{1}{p_{2}\left( b,i\right)}\right)\right)
\end{array}$



This gives us:



\textbf{Corollary 8: }

$\displaystyle  \begin{array}{l}
var( est2( C_{1} ;C_{2} ;a_{1} ,...,a_{k} ;b_{1} ,...,b_{k}))\\
\rightarrow \frac{( \Delta z_{1} -\Delta z_{2})^{2}}{( c_{1} -c_{2})^{4}}\left(\frac{c_{2}^{2}}{N_{1}}\left(\sum _{i=1}^{k}\frac{1}{p_{1}( a,i)} +\sum _{i=1}^{k}\frac{1}{p_{1}( b,i)}\right) +\frac{c_{1}^{2}}{N_{2}}\left(\sum _{i=1}^{k}\frac{1}{p_{2}( a,i)} +\sum _{i=1}^{k}\frac{1}{p_{2}( b,i)}\right)\right)
\end{array}$ as $\displaystyle N_{1} ,N_{2}\rightarrow \infty $\qed 



Recall that when LLM A = LLM B, we have that $\displaystyle c_{j} =\frac{\Delta z_{j} +\Delta b}{T}$. This gives us:



\textbf{Corollary 9:} If LLM A = LLM B, $\displaystyle
var\left( est2\left( C_{1} ;C_{2} ;a_{1} ,...,a_{k} ;b_{1} ,...,b_{k}\right)\right)$\\
$\rightarrow \frac{T^{2}}{\left( \Delta z_{1} -\Delta z_{2}\right)^{2}}\left(\frac{\left( \Delta z_{2} +\Delta b\right)^{2}}{N_{1}}\left(\sum _{i=1}^{k}\frac{1}{p_{1}\left( a,i\right)} +\sum _{i=1}^{k}\frac{1}{p_{1}\left( b,i\right)}\right) +\frac{\left( \Delta z_{1} +\Delta b\right)^{2}}{N_{2}}\left(\sum _{i=1}^{k}\frac{1}{p_{2}\left( a,i\right)} +\sum _{i=1}^{k}\frac{1}{p_{2}\left( b,i\right)}\right)\right)$



Proof:

$\displaystyle  \begin{array}{l}
\frac{\left( \Delta z_{1} -\Delta z_{2}\right)^{2}}{\left( c_{1} -c_{2}\right)^{4}}\left(\frac{c_{2}^{2}}{N_{1}}\left(\sum _{i=1}^{k}\frac{1}{p_{1}\left( a,i\right)} +\sum _{i=1}^{k}\frac{1}{p_{1}\left( b,i\right)}\right) +\frac{c_{1}^{2}}{N_{2}}\left(\sum _{i=1}^{k}\frac{1}{p_{2}\left( a,i\right)} +\sum _{i=1}^{k}\frac{1}{p_{2}\left( b,i\right)}\right)\right)\\
=\frac{\left( \Delta z_{1} -\Delta z_{2}\right)^{2}}{\left(\frac{\Delta z_{1} +\Delta b}{T} -\frac{\Delta z_{2} +\Delta b}{T}\right)^{4}}\left(\frac{\left(\frac{\Delta z_{2} +\Delta b}{T}\right)^{2}}{N_{1}}\left(\sum _{i=1}^{k}\frac{1}{p_{1}\left( a,i\right)} +\sum _{i=1}^{k}\frac{1}{p_{1}\left( b,i\right)}\right) +\frac{\left(\frac{\Delta z_{1} +\Delta b}{T}\right)^{2}}{N_{2}}\left(\sum _{i=1}^{k}\frac{1}{p_{2}\left( a,i\right)} +\sum _{i=1}^{k}\frac{1}{p_{2}\left( b,i\right)}\right)\right)\\
=\frac{T^{2}}{\left( \Delta z_{1} -\Delta z_{2}\right)^{2}}\left(\frac{\left( \Delta z_{2} +\Delta b\right)^{2}}{N_{1}}\left(\sum _{i=1}^{k}\frac{1}{p_{1}\left( a,i\right)} +\sum _{i=1}^{k}\frac{1}{p_{1}\left( b,i\right)}\right) +\frac{\left( \Delta z_{1} +\Delta b\right)^{2}}{N_{2}}\left(\sum _{i=1}^{k}\frac{1}{p_{2}\left( a,i\right)} +\sum _{i=1}^{k}\frac{1}{p_{2}\left( b,i\right)}\right)\right)
\end{array}$

\qed 



\subsubsection*{Heuristic Analysis:}

Looking at the case of LLM A = LLM B, it is trickier to find a heuristic for producing candidates for $\displaystyle C_{1} ,C_{2}$, and $\displaystyle a_{1} ,...,a_{k} ,b_{1} ,...,b_{k}$ than for $\displaystyle est$. In particular, there are 2 simultaneously competing effects at play: reducing variance corresponds to increasing $\displaystyle \left( \Delta z_{1} -\Delta z_{2}\right)^{2}$, but reducing $\displaystyle \Delta z{_{2}}^{2}$ and $\displaystyle \Delta z_{1}^{2}$.



Despite it being trickier to produce a theoretically-informed heuristic for $\displaystyle est2$ compared to $\displaystyle est$, $\displaystyle est2$ has a notable practical advantage: unlike $\displaystyle est$, $\displaystyle est2$ often works as a stand-alone estimator for determining the ownership identity of LLM B in practice.



To see why, let us recall that $\displaystyle est2$ converges to $\displaystyle \Delta b$ when LLM A = LLM B. From the definition of $\displaystyle \Delta b$, it follows that $\displaystyle \Delta b=0$ whenever logit biases are 0 or uniformly applied to the tokens $\displaystyle a_{1} ,...,a_{k} ,b_{1} ,...,b_{k}$. This means that for $\displaystyle \Delta b\neq 0$ to occur, an adversary would have to not only go out of their way to apply heterogeneous logit biases to $\displaystyle a_{1} ,...,a_{k} ,b_{1} ,...,b_{k}$, but they would have likely had to guess which tokens we were going to use in creating the estimator $\displaystyle est2$ in the first place (this, of course, precludes cleverer adversaries who automate adding small, random logit biases to every token). Moreover, if $\displaystyle a_{1} ,...,a_{k} ,b_{1} ,...,b_{k}$ are valuable tokens for LLM performance (e.g. numerical), then applying logit biases to them can significantly degrade the LLM's performance. In practice, logit biases are rarely applied, and when they are, they are often uniformly applied. This means that it is usually the case that $\displaystyle \Delta b=0$. Thus, in many practical scenarios, a simple test for whether LLM A = LLM B is given by checking if $\displaystyle est2$ converges to 0.



By contrast, $\displaystyle est$ estimates temperature $\displaystyle T$. There is no universal standard for what the default temperature is in practical LLM deployments. In practice, $\displaystyle T$ is often set anywhere from 0.7 to 1 and sometimes out of this range. One can make the counterargument that $\displaystyle T$ is often set to a multiple of 0.1 and therefore $\displaystyle est$'s convergence to a multiple of 0.1 is a sign that LLM A = LLM B. This is valid when this assumption is true, but not when the adversary simply sets the temperature of LLM B to an irregular value.



\subsubsection{Naive Method I}
\begin{mdframed}
    For the rest of this "Naive Method I" subsection, we are assuming all assumptions made during its presentation in the "Determining LLM Ownership Identity" section, in addition to the prevailing assumptions being made throughout the entire "Asymptotic Analysis: Independent $\displaystyle g_{j}$'s" section.
\end{mdframed}

With the assumptions and notation of Naive Method I, we have: $\displaystyle \phi \left( g_{1} ,g_{2} ,g_{3} ,g_{4}\right) =\frac{\Delta z_{1} -\Delta z_{2}}{g_{1} -g_{2}} -\frac{\Delta z_{3} -\Delta z_{4}}{g_{3} -g_{4}}$

Then $\displaystyle var\left( \phi \left( g_{1} ,g_{2} ,g_{3} ,g_{4}\right)\right) =var\left(\frac{\Delta z_{1} -\Delta z_{2}}{g_{1} -g_{2}}\right) +var\left(\frac{\Delta z_{3} -\Delta z_{4}}{g_{3} -g_{4}}\right)$ due to the independence of the $\displaystyle g_{j}$'s.



We can rewrite this as: $\displaystyle var\left( \phi \left( g_{1} ,g_{2} ,g_{3} ,g_{4}\right)\right) =var\left( est\left( C_{1} ;C_{2} ;a_{1} ,...,a_{k} ;b_{1} ,...,b_{k}\right)\right) +var\left( est\left( C_{3} ;C_{4} ;a_{1} ',...,a_{k'} ';b_{1} ',...,b_{k'} '\right)\right)$



Heuristic analysis of $\displaystyle var\left( \phi \left( g_{1} ,g_{2} ,g_{3} ,g_{4}\right)\right)$ then follows from the "The estimator est" subsection.



\subsubsection{Naive Method II}
\begin{mdframed}
    For the rest of this "Naive Method I" subsection, we are assuming all assumptions made during its presentation in the "Determining LLM Ownership Identity" section, in addition to the prevailing assumptions being made throughout the entire "Asymptotic Analysis: Independent $\displaystyle g_{j}$'s" section.
\end{mdframed}
With the assumptions and notation of Naive Method II, we have: $\displaystyle \phi \left( g_{1} ,g_{2} ,g_{3} ,g_{4}\right) =\frac{g_{2} \Delta z_{1} -g_{1} \Delta z_{2}}{g_{1} -g_{2}} -\frac{g_{4} \Delta z_{3} -g_{3} \Delta z_{4}}{g_{3} -g_{4}}$.



Then $\displaystyle var\left( \phi \left( g_{1} ,g_{2} ,g_{3} ,g_{4}\right)\right) =var\left(\frac{g_{2} \Delta z_{1} -g_{1} \Delta z_{2}}{g_{1} -g_{2}}\right) +var\left(\frac{g_{4} \Delta z_{3} -g_{3} \Delta z_{4}}{g_{3} -g_{4}}\right)$ due to the independence of the $\displaystyle g_{j}$'s.



We can rewrite this as: $\displaystyle var\left( \phi \left( g_{1} ,g_{2} ,g_{3} ,g_{4}\right)\right) =var\left( est2\left( C_{1} ;C_{2} ;a_{1} ,...,a_{k} ;b_{1} ,...,b_{k}\right)\right) +var\left( est2\left( C_{3} ;C_{4} ;a_{1} ,...,a_{k} ;b_{1} ,...,b_{k}\right)\right)$



Heuristic analysis of $\displaystyle var\left( \phi \left( g_{1} ,g_{2} ,g_{3} ,g_{4}\right)\right)$ then follows from the "The estimator est2" subsection.


\subsection{Deferred Proofs ($\star\star$)}

\textbf{Definitions:}
\begin{itemize}
\item Let $\displaystyle X_{j,i} =\frac{n_{j}\left( a,i\right)}{N_{j}}$ for $\displaystyle i=1:k_{j}$
\item Let $\displaystyle Y_{j,i} =\frac{n_{j}\left( b,i\right)}{N_{j}}$ for $\displaystyle i=1:k_{j}$
\end{itemize}

\textbf{Theorem 9:} $\displaystyle var\left( g_{j}\right)\rightarrow \frac{1}{N_{j}}\left(\sum _{i=1}^{k_{j}}\frac{1}{p_{j}\left( a,i\right)} +\sum _{i=1}^{k_{j}}\frac{1}{p_{j}\left( b,i\right)}\right)$ as $\displaystyle N_{j}\rightarrow \infty $

Proof:

Suppose that $\displaystyle N_{j}\rightarrow \infty $.

First of all, note that we have:

$\displaystyle g_{j}\left( X_{j,1} ,...,X_{j,k_{j}} ,Y_{j,1} ,...,Y_{j,k_{j}}\right) =\sum _{i=1}^{k_{j}} ln\left( N_{j} X_{j,i}\right) -\sum _{i=1}^{k_{j}} ln\left( N_{j} Y_{j,i}\right) =\sum _{i=1}^{k_{j}} ln\left( X_{j,i}\right) -\sum _{i=1}^{k_{j}} ln\left( Y_{j,i}\right)$

Note that $\displaystyle g_{j}\left( X_{j,1} ,...,X_{j,k_{j}} ,Y_{j,1} ,...,Y_{j,k_{j}}\right)$'s function is independent of $\displaystyle N_{j}$. Furthermore, note that $\displaystyle X_{j,i}$ and $\displaystyle Y_{j,i}$ are consistent estimators of $\displaystyle p_{j}\left( a,i\right)$ and $\displaystyle p_{j}\left( b,i\right)$ respectively. Therefore, by the delta method, we have:

$\displaystyle  \begin{array}{l}
g_{j}\left( X_{j,1} ,...,X_{j,k_{j}} ,Y_{j,1} ,...,Y_{j,k_{j}}\right) \approx g_{j}\left( p_{j}\left( a,1\right) ,...,p_{j}\left( a,k_{j}\right) ,p_{j}\left( b,1\right) ,...,p_{j}\left( b,k_{j}\right)\right)\\
+\sum _{i=1}^{k_{j}}\frac{\partial g}{\partial x_{j,i}}\left( p_{j}\left( a,1\right) ,...,p_{j}\left( a,k_{j}\right) ,p_{j}\left( b,1\right) ,...,p_{j}\left( b,k_{j}\right)\right)\left( X_{j,i} -p_{j}\left( a,i\right)\right)
\end{array}$

$\displaystyle +\sum _{i=1}^{k_{j}}\frac{\partial g}{\partial y_{j,i}}\left( p_{j}\left( a,1\right) ,...,p_{j}\left( a,k_{j}\right) ,p_{j}\left( b,1\right) ,...,p_{j}\left( b,k_{j}\right)\right)\left( Y_{j,i} -p_{j}\left( b,i\right)\right)$ (4)



Let $\displaystyle A_{i} =\frac{\partial g}{\partial x_{j,i}}\left( p_{j}\left( a,1\right) ,...,p_{j}\left( a,k_{j}\right) ,p_{j}\left( b,1\right) ,...,p_{j}\left( b,k_{j}\right)\right)$

Let $\displaystyle B_{i} =\frac{\partial g}{\partial y_{j,i}}\left( p_{j}\left( a,1\right) ,...,p_{j}\left( a,k_{j}\right) ,p_{j}\left( b,1\right) ,...,p_{j}\left( b,k_{j}\right)\right)$



Then we have:

$\displaystyle \begin{array}{l}
var( g_{j}( X_{j,1} ,...,X_{j,k_{j}} ,Y_{j,1} ,...,Y_{j,k_{j}}))\\
\approx var\left(\sum _{i=1}^{k_{j}} A_{i}( X_{j,i} -p_{j}( a,i)) +\sum _{i=1}^{k_{j}} B_{i}( Y_{j,i} -p_{j}( b,i))\right)\\
=\sum _{i=1}^{k_{j}} A_{i}^{2} var( X_{j,i} -p_{j}( a,i)) +\sum _{i=1}^{k_{j}} B_{i}^{2} var( Y_{j,i} -p_{j}( b,i))\\
+2\sum _{1\leqslant i< l\leqslant k_{j}} A_{i} A_{l} Cov( X_{j,i} -p_{j}( a,i) ,X_{j,l} -p_{j}( a,l))\\
+2\sum _{1\leqslant i< l\leqslant k_{j}} B_{i} B_{l} Cov( Y_{j,i} -p_{j}( b,i) ,Y_{j,l} -p_{j}( b,l))\\
+2\sum _{i=1}^{k_{j}}\sum _{l=1}^{k_{j}} A_{i} B_{l} Cov( X_{j,i} -p_{j}( a,i) ,Y_{j,l} -p_{j}( b,l))\\
=\sum _{i=1}^{k_{j}} A_{i}^{2} var( X_{j,i}) +\sum _{i=1}^{k_{j}} B_{i}^{2} var( Y_{j,i})\\
+2\sum _{1\leqslant i< l\leqslant k_{j}} A_{i} A_{l} Cov( X_{j,i} ,X_{j,l})\\
+2\sum _{1\leqslant i< l\leqslant k_{j}} B_{i} B_{l} Cov( Y_{j,i} ,Y_{j,l})\\
+2\sum _{i=1}^{k_{j}}\sum _{l=1}^{k_{j}} A_{i} B_{l} Cov( X_{j,i} ,Y_{j,l})\\
=\sum _{i=1}^{k_{j}} A_{i}^{2} var\left(\frac{n_{j}( a,i)}{N_{j}}\right) +\sum _{i=1}^{k_{j}} B_{i}^{2} var\left(\frac{n_{j}( b,i)}{N_{j}}\right)\\
+2\sum _{1\leqslant i< l\leqslant k_{j}} A_{i} A_{l} Cov\left(\frac{n_{j}( a,i)}{N_{j}} ,\frac{n_{j}( a,l)}{N_{j}}\right)\\
+2\sum _{1\leqslant i< l\leqslant k_{j}} B_{i} B_{l} Cov\left(\frac{n_{j}( b,i)}{N_{j}} ,\frac{n_{j}( b,l)}{N_{j}}\right)\\
+2\sum _{i=1}^{k_{j}}\sum _{l=1}^{k_{j}} A_{i} B_{l} Cov\left(\frac{n_{j}( a,i)}{N_{j}} ,\frac{n_{j}( b,l)}{N_{j}}\right)
\end{array}$

Due to (*) preventing overlapping generations among $\displaystyle a_{1,j} ,...,a_{k_{j} ,j} ,b_{1,j} ,...,b_{k_{j} ,j}$ for each $\displaystyle j$, the $\displaystyle n_{j}( a,i)$'s and $\displaystyle n_{j}( b,i)$'s are components of a multinomial distribution for each $\displaystyle j$. The following are facts from multinomial distributions:
\begin{itemize}
\item $\displaystyle var( n_{j}( a,i)) =N_{j} p_{j}( a,i)( 1-p_{j}( a,i))$
\item $\displaystyle var( n_{j}( b,i)) =N_{j} p_{j}( b,i)( 1-p_{j}( b,i))$
\item $\displaystyle Cov( n_{j}( a,i) ,n_{j}( a,l)) =-N_{j} p_{j}( a,i) p_{j}( a,l)$
\item $\displaystyle Cov( n_{j}( b,i) ,n_{j}( b,l)) =-N_{j} p_{j}( b,i) p_{j}( b,l)$
\item $\displaystyle Cov( n_{j}( a,i) ,n_{j}( b,l)) =-N_{j} p_{j}( a,i) p_{j}( b,l)$
\end{itemize}

So the above expression becomes:

$\displaystyle  \begin{array}{l}
\sum _{i=1}^{k_{j}}\frac{A_{i}^{2}}{N_{j}^{2}} N_{j} p_{j}( a,i)( 1-p_{j}( a,i)) +\sum _{i=1}^{k_{j}}\frac{B_{i}^{2}}{N_{j}^{2}} N_{j} p_{j}( b,i)( 1-p_{j}( b,i))\\
+2\sum _{1\leqslant i< l\leqslant k_{j}}\frac{A_{i} A_{l}}{N_{j}^{2}}( -N_{j} p_{j}( a,i) p_{j}( a,l))\\
+2\sum _{1\leqslant i< l\leqslant k_{j}}\frac{B_{i} B_{l}}{N_{j}^{2}}( -N_{j} p_{j}( b,i) p_{j}( b,l))\\
+2\sum _{i=1}^{k_{j}}\sum _{l=1}^{k_{j}}\frac{A_{i} B_{l}}{N_{j}^{2}}( -N_{j} p_{j}( a,i) p_{j}( b,l))\\
=\frac{1}{N_{j}}(\sum _{i=1}^{k_{j}} A_{i}^{2} p_{j}( a,i)( 1-p_{j}( a,i)) +\sum _{i=1}^{k_{j}} B_{i}^{2} p_{j}( b,i)( 1-p_{j}( b,i))\\
-2\sum _{1\leqslant i< l\leqslant k_{j}} A_{i} A_{l}( p_{j}( a,i) p_{j}( a,l))\\
-2\sum _{1\leqslant i< l\leqslant k_{j}} B_{i} B_{l}( p_{j}( b,i) p_{j}( b,l))\\
-2\sum _{i=1}^{k_{j}}\sum _{l=1}^{k_{j}} A_{i} B_{l}( p_{j}( a,i) p_{j}( b,l)))
\end{array}$



We now compute $\displaystyle A_{i}$ and $\displaystyle B_{i}$.

$\displaystyle g_{j}( x_{j,1} ,...,x_{j,k_{j}} ,y_{j,1} ,...,y_{j,k_{j}}) =\sum _{i=1}^{k_{j}} ln( X_{j,i}) -\sum _{i=1}^{k_{j}} ln( Y_{j,i})$

So $\displaystyle \frac{\partial g}{\partial x_{j,i}}( x_{j,1} ,...,x_{j,k_{j}} ,y_{j,1} ,...,y_{j,k_{j}}) =\frac{1}{x_{j,i}}$ and $\displaystyle \frac{\partial g}{\partial y_{j,i}}( x_{j,1} ,...,x_{j,k_{j}} ,y_{j,1} ,...,y_{j,k_{j}}) =-\frac{1}{y_{j,i}}$



Then we have:

$\displaystyle  \begin{array}{l}
A_{i} =\frac{\partial g}{\partial x_{j,i}}( p_{j}( a,1) ,...,p_{j}( a,k_{j}) ,p_{j}( b,1) ,...,p_{j}( b,k_{j}))\\
=\frac{1}{p_{j}( a,i)}
\end{array}$

And similarly:

$\displaystyle \begin{array}{l}
B_{i} =\frac{\partial g}{\partial y_{j,i}}( p_{j}( a,1) ,...,p_{j}( a,k_{j}) ,p_{j}( b,1) ,...,p_{j}( b,k_{j}))\\
=-\frac{1}{p_{j}( b,i)}
\end{array}$



Therefore, the above expression becomes:

$\displaystyle  \begin{array}{l}
\frac{1}{N_{j}}(\sum _{i=1}^{k_{j}}\frac{1}{p_{j}( a,i)^{2}} p_{j}( a,i)( 1-p_{j}( a,i)) +\sum _{i=1}^{k_{j}}\left( -\frac{1}{p_{j}( b,i)}\right)^{2} p_{j}( b,i)( 1-p_{j}( b,i))\\
-2\sum _{1\leqslant i< l\leqslant k_{j}}\frac{1}{p_{j}( a,i)}\frac{1}{p_{j}( a,l)}( p_{j}( a,i) p_{j}( a,l))\\
-2\sum _{1\leqslant i< l\leqslant k_{j}}\left( -\frac{1}{p_{j}( b,i)}\right)\left( -\frac{1}{p_{j}( b,l)}\right)( p_{j}( b,i) p_{j}( b,l))\\
-2\sum _{i=1}^{k_{j}}\sum _{l=1}^{k_{j}}\frac{1}{p_{j}( a,i)}\left( -\frac{1}{p_{j}( b,l)}\right)( p_{j}( a,i) p_{j}( b,l)))\\
=\frac{1}{N_{j}}(\sum _{i=1}^{k_{j}}\left(\frac{1}{p_{j}( a,i)} -1\right) +\sum _{i=1}^{k_{j}}\left(\frac{1}{p_{j}( b,i)} -1\right)\\
-2\sum _{1\leqslant i< l\leqslant k_{j}} 1\\
-2\sum _{1\leqslant i< l\leqslant k_{j}} 1\\
+2\sum _{i=1}^{k_{j}}\sum _{l=1}^{k_{j}} 1)\\
=\frac{1}{N_{j}}\left(\sum _{i=1}^{k_{j}}\frac{1}{p_{j}( a,i)} +\sum _{i=1}^{k_{j}}\frac{1}{p_{j}( b,i)}\right)
\end{array}$

\qed 



\textbf{Theorem 10:} Each $\displaystyle g_{j}$ is asymptotically a normal random variable as $\displaystyle N_{j}\rightarrow \infty $.

Proof:

Suppose $\displaystyle N_{j}\rightarrow \infty $.

From equation (4), we have the following:

$\displaystyle g_{j}\left( X_{j,1} ,...,X_{j,k_{j}} ,Y_{j,1} ,...,Y_{j,k_{j}}\right) \approx c+\sum _{i=1}^{k_{j}} d_{i}\left( X_{j,i} -p_{j}\left( a,i\right)\right) +\sum _{i=1}^{k_{j}} e_{i}\left( Y_{j,i} -p_{j}\left( b,i\right)\right)$

For some constants $\displaystyle c,d_{i} ,e_{i}$. Since normality is preserved under an additive constant, it suffices to show that the following is normal:

$\displaystyle  \begin{array}{l}
\sum _{i=1}^{k_{j}} d_{i} X_{j,i} +\sum _{i=1}^{k_{j}} e_{i} Y_{j,i}\\
=\sum _{i=1}^{k_{j}} d_{i}\frac{n_{j}\left( a,i\right)}{N_{j}} +\sum _{i=1}^{k_{j}} e_{i}\frac{n_{j}\left( b,i\right)}{N_{j}}
\end{array}$

Write $\displaystyle n_{j}\left( a,i\right)$ as a sum of IID indicator random variables $\displaystyle I_{j}\left( a,i,t\right)$ for $\displaystyle t=1:N_{j}$, each indicator random variable indicating whether $\displaystyle a_{i}$ was a prefix string of the $\displaystyle t$th generated output of LLM B on context $\displaystyle C_{j}$. Define $\displaystyle I_{j}\left( b,i,t\right)$ similarly.

Thus, we have:

$\displaystyle \begin{array}{l}
\sum _{i=1}^{k_{j}} d_{i}\frac{n_{j}\left( a,i\right)}{N_{j}} +\sum _{i=1}^{k_{j}} e_{i}\frac{n_{j}\left( b,i\right)}{N_{j}}\\
=\sum _{i=1}^{k_{j}} d_{i}\frac{\sum _{t=1}^{N_{j}} I_{j}\left( a,i,t\right)}{N_{j}} +\sum _{i=1}^{k_{j}} e_{i}\frac{\sum _{t=1}^{N_{j}} I_{j}\left( b,i,t\right)}{N_{j}}\\
=\frac{1}{N_{j}}\sum _{t=1}^{N_{j}}\left(\sum _{i=1}^{k_{j}} d_{i} I_{j}\left( a,i,t\right) +\sum _{i=1}^{k_{j}} e_{i} I_{j}\left( b,i,t\right)\right)
\end{array}$

Take $\displaystyle Z_{t} =\sum _{i=1}^{k_{j}} d_{i} I_{j}\left( a,i,t\right) +\sum _{i=1}^{k_{j}} e_{i} I_{j}\left( b,i,t\right)$. Note that the $\displaystyle Z_{t}$'s are IID.

The expression then becomes $\displaystyle \frac{1}{N_{j}}\sum _{t=1}^{N_{j}} Z_{t}$, which is asymptotically normal by the central limit theorem.\qed 



\textbf{Lemma 4: }$\displaystyle Bias\left( g_{j}\right) =E\left[ g_{j}\right] -c_{j} =O\left( 1/N_{j}\right)$

Proof:

Recall from the proof of theorem 9 that we may write $\displaystyle g_{j} =\sum _{i=1}^{k_{j}} ln\left( X_{j,i}\right) -\sum _{i=1}^{k_{j}} ln\left( Y_{j,i}\right)$.

Then:

$\displaystyle  \begin{array}{l}
E\left[ g_{j}\right] -c_{j} =\sum _{i=1}^{k_{j}} E\left[ ln\left( X_{j,i}\right)\right] -\sum _{i=1}^{k_{j}} E\left[ ln\left( Y_{j,i}\right)\right] -\left(\sum _{i=1}^{k_{j}} ln\left( p_{j}\left( a,i\right)\right) -\sum _{i=1}^{k_{j}} ln\left( p_{j}\left( b,i\right)\right)\right)\\
=\sum _{i=1}^{k_{j}}\left( E\left[ ln\left( X_{j,i}\right)\right] -ln\left( p_{j}\left( a,i\right)\right)\right) -\sum _{i=1}^{k_{j}}\left( E\left[ ln\left( Y_{j,i}\right)\right] -ln\left( p_{j}\left( b,i\right)\right)\right)
\end{array}$

Since $\displaystyle k_{j}$'s are fixed, it suffices to show that each $\displaystyle E\left[ ln\left( X_{j,i}\right)\right] -ln\left( p_{j}\left( a,i\right)\right)$ and each $\displaystyle E\left[ ln\left( Y_{j,i}\right)\right] -ln\left( p_{j}\left( b,i\right)\right)$ term is $\displaystyle O\left( 1/N_{j}\right)$. We will show this for $\displaystyle E\left[ ln\left( X_{j,i}\right)\right] -ln\left( p_{j}\left( a,i\right)\right)$. The proof for the $\displaystyle E\left[ ln\left( Y_{j,i}\right)\right] -ln\left( p_{j}\left( b,i\right)\right)$ case is identical.

We have by the 2nd-order Taylor approximation that:

$\displaystyle ln\left( X_{j,i}\right) =ln\left( p_{j}\left( a,i\right)\right) +\frac{1}{p_{j}\left( a,i\right)}\left( X_{j,i} -p_{j}\left( a,i\right)\right) -\frac{1}{2p_{j}\left( a,i\right)^{2}}\left( X_{j,i} -p_{j}\left( a,i\right)\right)^{2} +R$ where $\displaystyle R$ is the remainder term.

Thus, we have:

$\displaystyle E\left[ ln\left( X_{j,i}\right)\right] -ln\left( p_{j}\left( a,i\right)\right) =\frac{1}{p_{j}\left( a,i\right)} E\left[\left( X_{j,i} -p_{j}\left( a,i\right)\right)\right] -\frac{1}{2p_{j}\left( a,i\right)^{2}} E\left[\left( X_{j,i} -p_{j}\left( a,i\right)\right)^{2}\right] +E\left[ R\right]$

Here are some standard facts from binomial distributions:
\begin{itemize}
\item $\displaystyle E\left[\left( X_{j,i} -p_{j}\left( a,i\right)\right)^{2}\right] =O\left( 1/N_{j}\right)$
\item $\displaystyle |X_{j,i} -p_{j}\left( a,i\right) |=O_{p}\left( 1/N_{k}^{1/2}\right)$
\end{itemize}

Note that the $\displaystyle \frac{1}{p_{j}\left( a,i\right)} E\left[\left( X_{j,i} -p_{j}\left( a,i\right)\right)\right]$ term is simply 0.

Furthermore, the $\displaystyle -\frac{1}{2p_{j}\left( a,i\right)^{2}} E\left[\left( X_{j,i} -p_{j}\left( a,i\right)\right)^{2}\right]$ term is $\displaystyle O\left( 1/N_{j}\right)$.

$\displaystyle |R|=O\left( |X_{j,i} -p_{j}\left( a,i\right) |^{3}\right)$ follows from a standard bound for Taylor polynomial remainders.

Thus, $\displaystyle |E\left[ R\right] |\leqslant E\left[ |R|\right] =O\left( E\left[ |X_{j,i} -p_{j}\left( a,i\right) |^{3}\right]\right)$.

Then $\displaystyle |X_{j,i} -p_{j}\left( a,i\right) |^{3} =O_{p}\left( 1/N_{k}^{3/2}\right) \Longrightarrow O\left( E\left[ |X_{j,i} -p_{j}\left( a,i\right) |^{3}\right]\right) =O\left( 1/N_{k}^{3/2}\right) =O\left( 1/N_{k}\right)$.

Therefore, $\displaystyle |E\left[ R\right] |=O\left( 1/N_{k}\right)$.

Therefore, $\displaystyle E\left[ ln\left( X_{j,i}\right)\right] -ln\left( p_{j}\left( a,i\right)\right) =O\left( 1/N_{k}\right)$.

This completes the proof.\qed 

\textbf{Theorem 14: }$\displaystyle Bias\left( \phi \left( g_{1} ,...,g_{m}\right)\right) =E\left[ \phi \left( g_{1} ,...,g_{m}\right)\right] -\phi \left( c_{1} ,...,c_{m}\right) =O\left(\sum _{j=1}^{m} 1/N_{j}\right)$

Proof:

Let $\displaystyle g=\left( g_{1} ,...,g_{m}\right)^{T}$ and $\displaystyle c=\left( c_{1} ,...,c_{m}\right)^{T}$. By the 2nd-order Taylor approximation, we have that:

$\displaystyle \phi \left( g\right) =\phi \left( c\right) +\nabla \phi \left( c\right)^{T}\left( g-c\right) +\frac{1}{2}\left( g-c\right)^{T} H_{\phi }\left( c\right)\left( g-c\right)^{T} +R$.

Thus, we have:

$\displaystyle E\left[ \phi \left( g\right)\right] -\phi \left( c\right) =\nabla \phi \left( c\right)^{T} E\left[ g-c\right] +E\left[\frac{1}{2}\left( g-c\right)^{T} H_{\phi }\left( c\right)\left( g-c\right)^{T}\right] +E\left[ R\right]$

We have $\displaystyle \nabla \phi \left( c\right)^{T} E\left[ g-c\right] =O\left(\sum _{j=1}^{m} 1/N_{j}\right)$ by lemma 4.

We have:

$\displaystyle  \begin{array}{l}
|E\left[\frac{1}{2}\left( g-c\right)^{T} H_{\phi }\left( c\right)\left( g-c\right)^{T}\right] |\leqslant E\left[ |\frac{1}{2}\left( g-c\right)^{T} H_{\phi }\left( c\right)\left( g-c\right) |\right]\\
\leqslant E\left[\frac{1}{2} ||H_{\phi }\left( c\right) ||_{2} ||g-c||_{2}^{2}\right]\\
=\frac{1}{2} ||H_{\phi }\left( c\right) ||_{2} E\left[ ||g-c||_{2}^{2}\right]\\
=const\cdot \sum _{j=1}^{m} E\left[\left( g_{j} -c_{j}\right)^{2}\right]
\end{array}$

We then have:

$\displaystyle E\left[\left( g_{j} -c_{j}\right)^{2}\right] =var\left( g_{j} -c_{j}\right) +E\left[ g_{j} -c_{j}\right]^{2}$

$\displaystyle =var\left( g_{j}\right) +O\left( 1/N_{j}\right)^{2}$ (by lemma 4)

$\displaystyle =O\left( 1/N_{j}\right) +O\left( 1/N_{j}\right)^{2}$ (by theorem 9)

$\displaystyle =O\left( 1/N_{j}\right)$

Thus, $\displaystyle |E\left[\frac{1}{2}\left( g-c\right)^{T} H_{\phi }\left( c\right)\left( g-c\right)^{T}\right] |=O\left(\sum _{j=1}^{m} 1/N_{j}\right)$.

$\displaystyle |R|=O\left( ||g-c||^{3}\right)$ follows from a standard bound for Taylor polynomial remainders.

Then:

$\displaystyle \begin{array}{l}
|E\left[ R\right] |\leqslant E\left[ |R|\right]\\
=E\left[ O\left( ||g-c||^{3}\right)\right]\\
\leqslant E\left[ O\left(\left(\sum _{j=1}^{m} |g_{j} -c_{j} |\right)^{3}\right)\right]
\end{array}$

Claim that we'll prove at the end of this proof: $\displaystyle |g_{j} -c_{j} |=O_{p}\left( 1/\sqrt{N_{j}}\right)$

Then:

$\displaystyle  \begin{array}{l}
E\left[ O\left(\left(\sum _{j=1}^{m} |g_{j} -c_{j} |\right)^{3}\right)\right] =O\left(\left(\sum _{j=1}^{m} 1/\sqrt{N_{j}}\right)^{3}\right)\\
=O\left(\left( 1/\sqrt{N_{min}}\right)^{3}\right)\\
=O\left( 1/N_{min}^{-3/2}\right)\\
=O\left(\sum _{j=1}^{m} 1/N_{j}^{-3/2}\right)\\
=O\left(\sum _{j=1}^{m} 1/N_{j}\right)
\end{array}$

Thus, $\displaystyle |E\left[ R\right] |=O\left(\sum _{j=1}^{m} 1/N_{j}\right)$. We conclude that $\displaystyle E\left[ \phi \left( g\right)\right] -\phi \left( c\right) =O\left(\sum _{j=1}^{m} 1/N_{j}\right)$.

It now remains to prove $\displaystyle |g_{j} -c_{j} |=O_{p}\left( 1/\sqrt{N_{j}}\right)$.

By bias-variance decomposition, we have:

$\displaystyle E\left[\left( g_{j} -c_{j}\right)^{2}\right] =var\left( g_{j}\right) +\left( E\left[ g_{j}\right] -c_{j}\right)^{2}$

$\displaystyle =O\left( 1/N_{j}\right)$ by theorem 9 and lemma 4.

This implies that $\displaystyle E\left[ N_{j}\left( g_{j} -c_{j}\right)^{2}\right] =O\left( 1\right)$.

Markov's inequality gives us:

$\displaystyle P\left( |g_{j} -c_{j} | >\frac{M}{\sqrt{N}}\right) =P\left( N_{j}\left( g_{j} -c_{j}\right)^{2}  >M^{2}\right) \leqslant \frac{E\left[ N_{j}\left( g_{j} -c_{j}\right)^{2}\right]}{M^{2}} \leqslant \frac{const}{M^{2}}$.

This implies $\displaystyle \ |g_{j} -c_{j} |=O_{p}\left( 1/\sqrt{N_{j}}\right)$.\qed 



\section{Asymptotic Analysis: Non-independent $g_{j}$'s (The General Case)}
\begin{mdframed}
    \begin{itemize}
    \item Throughout this section, $\phi(x_1, \dots, x_m)$ will be an arbitrary analytic function unless otherwise stated.
    \item Throughout this section, we will be assuming $\dagger\dagger$ instead of $\dagger$ (which is a special case of $\dagger\dagger$; $\dagger\dagger$ is a relaxation of $\dagger$).
\end{itemize}
\end{mdframed}


The "Asymptotic Analysis: Independent $\displaystyle g_{j}$'s" section relied on the assumption that the $\displaystyle g_{j}$'s were decoupled and independent. In other words, each $\displaystyle g_{j}$ was assumed to be constructed from an independent, self-contained experiment involving running LLM B many times on context $\displaystyle C_{j}$ and recording number of occurrences of $\displaystyle a_{1,j} ,...,a_{k_{j} ,j} ,b_{1,j} ,...,b_{k_{j} ,j}$. The "What is an experiment?" subsection provided an illustrative example for a scenario where different $\displaystyle g_{j}$'s can become coupled and non-independent. We now consider the general case where independence isn't assumed in this section.

The previous section used nice properties such as independence and mutually exclusive events to derive elegant exact formulae, a luxury we do not have here. Indeed, if we chased exact formulae like in the previous section, we would end up with something hairy quite quickly due to our limited structural assumptions. Thus, we focus on high-level structure in this section - our analyses will be framed in terms of big-$O$ and big-$\Theta$. We use the previous section's results as guideposts and attempt to derive theorems that are \textit{structurally} analogous to them. Fortunately, many analogous theorems can be proven in the non-independent case that mirror the structure of the theorems in the independent case.

We will present these theorems in the language of \textit{experiments}, which was first introduced in "What is an experiment?"


\subsection{The Core Theorems (General Case)}



\textbf{Theorem 16: }$\displaystyle \phi \left( g_{1} ,...,g_{m}\right)$ is asymptotically a normal random variable as $\displaystyle N^{\left( 1\right)} ,...,N^{\left( w\right)}\rightarrow \infty $.



Proof:

First, note that by the law of large numbers, for any $\displaystyle n_{l,r} \in O_{r}$, $\displaystyle \frac{n_{l,r}}{N^{\left( r\right)}}\rightarrow p_{l,r}$ for some probability $\displaystyle p_{l,r}$ as $\displaystyle N^{\left( r\right)}\rightarrow \infty $. Define $\displaystyle X_{l}^{\left( r\right)} =\frac{n_{l,r}}{N^{\left( r\right)}}$.



Thus far, we've been treating $\displaystyle g_{j}$'s as the atomic building blocks of $\displaystyle \phi \left( g_{1} ,...,g_{m}\right)$, but there is an object that is even more granular: $\displaystyle n_{j}\left( a,i\right)$'s and $\displaystyle n_{j}\left( b,i\right)$'s. Recall the definition that $\displaystyle g_{j} =ln\left(\prod _{i=1}^{k_{j}}\frac{n_{j}\left( a,i\right)}{n_{j}\left( b,i\right)}\right) =\sum _{i=1}^{k_{j}} ln\left( n_{j}\left( a,i\right)\right) -\sum _{i=1}^{k_{j}} ln\left( n_{j}\left( b,i\right)\right)$. Thus, we may view $\displaystyle \phi \left( g_{1} ,...,g_{m}\right)$ as a function of $\displaystyle n_{j}\left( a,i\right)$'s and $\displaystyle n_{j}\left( b,i\right)$'s instead of as a function of $\displaystyle g_{j}$'s. Equivalently, we may view $\displaystyle \phi \left( g_{1} ,...,g_{m}\right)$ as a function of $\displaystyle n_{l,r}$'s. Equivalently, we may view $\displaystyle \phi \left( g_{1} ,...,g_{m}\right)$ as a function of $\displaystyle X_{l}^{\left( r\right)}$'s. 



We wish to apply the delta method on $\displaystyle \phi \left( g_{1} ,...,g_{m}\right)$ as a function of $\displaystyle X_{l}^{\left( r\right)}$'s, but we must first ensure that when $\displaystyle \phi \left( g_{1} ,...,g_{m}\right)$ is expressed as a function of $\displaystyle X_{l}^{\left( r\right)}$'s, it is a fixed function that does not depend on the $\displaystyle N^{\left( r\right)}$'s.



For every $\displaystyle j$, we know that $\displaystyle n_{j}\left( a,i\right)$ and $\displaystyle n_{j}\left( b,i\right)$ for $\displaystyle i=1:k_{j}$ all belong to the same $\displaystyle O_{r}$. Fix $\displaystyle r$. Consider again the equation $\displaystyle g_{j} =\sum _{i=1}^{k_{j}} ln\left( n_{j}\left( a,i\right)\right) -ln\left( n_{j}\left( b,i\right)\right)$. Let us consider index $\displaystyle i$ in this summation. Refer to $\displaystyle n_{j}\left( a,i\right)$ and $\displaystyle n_{j}\left( b,i\right)$ as $\displaystyle n_{l_{1} ,r}$ and $\displaystyle n_{l_{2} ,r}$ respectively. Then $\displaystyle ln\left( n_{j}\left( a,i\right)\right) -ln\left( n_{j}\left( b,i\right)\right) =ln\left( N^{\left( r\right)} X_{l_{1}}^{\left( r\right)}\right) -ln\left( N^{\left( r\right)} X_{l_{2}}^{\left( r\right)}\right) =ln\left( X_{l_{1}}^{\left( r\right)}\right) -ln\left( X_{l_{2}}^{\left( r\right)}\right)$. So we see that each term in the summation of $\displaystyle g_{j}$ does not depend on $\displaystyle N^{\left( r\right)}$. Therefore, $\displaystyle g_{j}$ is a fixed function of $\displaystyle X_{l}^{\left( r\right)}$ terms. Thus, $\displaystyle \phi \left( g_{1} ,...,g_{m}\right)$ is a fixed function of $\displaystyle X_{l}^{\left( r\right)}$ terms.



We may now apply the delta method:

$\displaystyle \phi \left( g_{1} ,...,g_{m}\right) \approx c+\sum _{r=1}^{w}\sum _{l=1}^{k^{\left( r\right)}} d_{l,r}\left( X_{l}^{\left( r\right)} -p_{l,r}\right)$ as $\displaystyle N^{\left( 1\right)} ,...,N^{\left( w\right)}\rightarrow \infty $ for some fixed constants $\displaystyle c$ and $\displaystyle d_{l,r} ,r=1:w,l=1:k^{\left( r\right)}$.



Note that each $\displaystyle \sum _{l=1}^{k^{\left( r\right)}} d_{l,r} X_{l}^{\left( r\right)}$ is independent for $\displaystyle r=1:w$ due to the $\displaystyle O_{r}$'s being statistically independent. Therefore, to prove asymptotic normality of $\displaystyle \phi \left( g_{1} ,...,g_{m}\right)$, it suffices to then prove that $\displaystyle \sum _{l=1}^{k^{\left( r\right)}} d_{l,r} X_{l}^{\left( r\right)}$ is asymptotically normal.



We have:

$\displaystyle \sum _{l=1}^{k^{\left( r\right)}} d_{l,r} X_{l}^{\left( r\right)} =\sum _{l=1}^{k^{\left( r\right)}} d_{l,r}\frac{n_{l,r}}{N^{\left( r\right)}}$

Write each $\displaystyle n_{l,r}$ as a sum of IID indicator random variables $\displaystyle I_{r}\left( l,t\right)$ for $\displaystyle t=1:N^{\left( r\right)}$, which record whether or not the event that $\displaystyle n_{l,r}$ is tracking occurred on trial/run $\displaystyle t$ of experiment $\displaystyle r$. We now have:

$\displaystyle  \begin{array}{l}
\sum _{l=1}^{k^{\left( r\right)}} d_{l,r}\frac{n_{l,r}}{N^{\left( r\right)}} =\sum _{l=1}^{k^{\left( r\right)}} d_{l,r}\frac{\sum _{t=1}^{N^{\left( r\right)}} I_{r}\left( l,t\right)}{N^{\left( r\right)}}\\
=\frac{1}{N^{\left( r\right)}}\sum _{t=1}^{N^{\left( r\right)}}\sum _{l=1}^{k^{\left( r\right)}} d_{l,r} I_{r}\left( l,t\right)
\end{array}$

Take $\displaystyle Z_{t}^{\left( r\right)} =\sum _{l=1}^{k^{\left( r\right)}} d_{l,r} I_{r}\left( l,t\right)$. Note that the $\displaystyle Z_{t}^{\left( r\right)}$'s are IID.

The expression then becomes $\displaystyle \frac{1}{N^{\left( r\right)}}\sum _{t=1}^{N^{\left( r\right)}} Z_{t}^{\left( r\right)}$, which is asymptotically normal by the central limit theorem. This completes the proof.\qed 



\textbf{Theorem 17: }$\displaystyle var\left( \phi \left( g_{1} ,...,g_{m}\right)\right)$ is $\displaystyle \Theta \left(\sum _{r=1}^{w}\frac{1}{N^{\left( r\right)}}\right)$ as $\displaystyle N^{\left( 1\right)} ,...,N^{\left( w\right)}\rightarrow \infty $

Proof:

From theorem 16's proof, we see that as $\displaystyle N^{\left( 1\right)} ,...,N^{\left( w\right)}\rightarrow \infty $, $\displaystyle \phi \left( g_{1} ,...,g_{m}\right)$ can be written, up to an additive constant, as $\displaystyle \sum _{r=1}^{w}\sum _{l=1}^{k^{\left( r\right)}} d_{l,r} X_{l}^{\left( r\right)} =\sum _{r=1}^{w}\left(\frac{1}{N^{\left( r\right)}}\sum _{t=1}^{N^{\left( r\right)}} Z_{t}^{\left( r\right)}\right)$.

Since different $\displaystyle O_{r}$'s are statistically independent, we have:

$\displaystyle var( \phi ( g_{1} ,...,g_{m})) =var\left(\sum _{r=1}^{w}\sum _{l=1}^{k^{( r)}} d_{l,r} X_{l}^{( r)}\right) =var\left(\sum _{r=1}^{w}\left(\frac{1}{N^{( r)}}\sum _{t=1}^{N^{( r)}} Z_{t}^{( r)}\right)\right) =\sum _{r=1}^{w} var\left(\frac{1}{N^{( r)}}\sum _{t=1}^{N^{( r)}} Z_{t}^{( r)}\right)$

Since the $\displaystyle Z_{t}^{\left( r\right)}$'s are IID for any given $\displaystyle r$, we can give $\displaystyle Z_{t}^{\left( r\right)} ,t=1:N^{\left( r\right)}$ the common variance $\displaystyle var^{\left( r\right)}$ and write:

$\displaystyle \sum _{r=1}^{w} var\left(\frac{1}{N^{\left( r\right)}}\sum _{t=1}^{N^{\left( r\right)}} Z_{t}^{\left( r\right)}\right) =\sum _{r=1}^{w}\frac{1}{N^{\left( r\right)}} var^{\left( r\right)} =\Theta \left(\sum _{r=1}^{w}\frac{1}{N^{\left( r\right)}}\right)$\qed 

Comment 1: Since we know $\displaystyle var\left( \phi \left( g_{1} ,...,g_{m}\right)\right)$ is $\displaystyle O\left(\sum _{r=1}^{w}\frac{1}{N^{\left( r\right)}}\right)$ as $\displaystyle N^{\left( 1\right)} ,...,N^{\left( w\right)}\rightarrow \infty $, then taking the $\displaystyle N^{\left( r\right)}$'s to be fixed proportions of $\displaystyle N$ (the total number of samples/calls/runs across all experiments) results in $\displaystyle O\left( 1/N\right)$ variance decay for $\displaystyle \phi \left( g_{1} ,...,g_{m}\right)$. This matches the best possible parametric decay rate $\displaystyle O\left( 1/N\right)$ for the variance of an estimator satisfying standard regularity conditions.



Comment 2: $\displaystyle N=N^{\left( 1\right)} +...+N^{\left( w\right)}$ is the number of calls to LLM B in the generic, non-independent $\displaystyle g_{j}$ setting. In general, $\displaystyle N_{1} +...+N_{m}  >N$ because the sum double-counts calls, since multiple $\displaystyle N_{j}$'s may come from the exact same experiment. $\displaystyle N=N_{1} +...+N_{m}$ is only true if we assume each $\displaystyle g_{j}$ was created from a self-contained, independent experiment (i.e. what was assumed in the previous section).



\textbf{Theorem 18:} $\displaystyle \phi \left( g_{1} ,...,g_{m}\right)$ is a consistent estimator of $\displaystyle \phi \left( c_{1} ,...,c_{m}\right)$ as $\displaystyle N^{\left( 1\right)} ,...,N^{\left( w\right)}\rightarrow \infty $.



Proof:

This is a restatement of Lemma 2. Note that at no intermediate step of Lemma 2's proof was the assumption of $\displaystyle g_{j}$ independence ever used, making it valid in this section as well.\qed 



\textbf{Theorem 19: }$\displaystyle Bias( \phi ( g_{1} ,...,g_{m})) =E[ \phi ( g_{1} ,...,g_{m})] -\phi ( c_{1} ,...,c_{m}) =O\left(\sum _{r=1}^{w} 1/N^{( r)}\right)$

Proof:

Theorem 14 from the previous section tells us that $\displaystyle Bias( \phi ( g_{1} ,...,g_{m})) =O\left(\sum _{j=1}^{m} 1/N_{j}\right)$. Note that we are allowed to apply theorem 14 in this setting since the proof of theorem 14 did not require assuming the $\displaystyle g_{j}$'s are independent.



We then have $\displaystyle O\left(\sum _{j=1}^{m} 1/N_{j}\right) =O\left( m\sum _{r=1}^{w} 1/N^{( r)}\right) =O\left(\sum _{r=1}^{w} 1/N^{( r)}\right)$, since $\displaystyle m$ is fixed.\qed 



Comment: From theorem 16, we have that $\displaystyle \sigma ( \phi ( g_{1} ,...,g_{m}))$ decays at rate $\displaystyle \Theta \left(\sqrt{\sum _{r=1}^{w} 1/N^{( r)}}\right)$, which asymptotically dominates $\displaystyle Bias( \phi ( g_{1} ,...,g_{m})) =O\left(\sum _{r=1}^{w} 1/N^{( r)}\right)$. In other words, $\displaystyle Bias( \phi ( g_{1} ,...,g_{m})) =o( \sigma ( \phi ( g_{1} ,...,g_{m}))$. Combined with the fact that $\displaystyle \phi ( g_{1} ,...,g_{m})$ is asymptotically normal, we may treat bias as negligible ($\displaystyle =0$) as $\displaystyle N^{( 1)} ,...,N^{( w)}\rightarrow \infty $ in many analyses.



\section{Conclusion}
\subsection{Summary}
In this paper, we analyzed the problem of detecting whether an anonymous LLM is a stolen version of a known one. The adversary this paper addresses is assumed to have perturbed the model via a set of hidden, standard token-distribution-modifying decoding parameters.

The central negative result proven in this paper is that such a set of perturbations is not enough to meaningfully conceal the LLM's identity; it was shown that a consistent, non-invasive, black-box test can indeed be constructed to determine its provenance.

In terms of efficiency, the tests developed in this paper depend on an estimator that converges in $O(1/N)$, matching the optimal convergence speed for canonical regular estimators. Furthermore, we conducted analyses to provide a broad theoretical framework for characterizing the tests' asymptotic behavior and improving their efficiency.

Along the way, we found several surprising intermediate results - for example, we showed that one can asymptotically recover the exact temperature and logit biases of tokens for LLM B when LLM A and LLM B are identical, further demonstrating that decoding parameters have significantly limited ability to destroy key information, even in a black-box setting.

\subsection{What this means for the proposed formal program}
As far as we have found, this paper is the first constructive evidence of the feasibility of the formal program whose aim is to classify threat scenarios which admit 'theoretically foolproof' (in the sense of decidability, statistical consistency, etc.) methods of LLM ownership identification. At the beginning of this paper, we reasoned that by necessity and practicality, decoding parameters are the first adversarial perturbation class that must be addressed by this program to provide credence for its feasibility. By tractably proving that an asymptotically strong test for it exists, this program may proceed with substantiated optimism in next steps to classify further threat scenarios.

\subsection{Future work}
In addition to carrying out the formal program, there is significant open room for future work that can be undertaken regarding the specific methods developed in this paper.

The methods developed here are well-positioned for empirical investigations conducting demos and experiments. In particular, it would be interesting to investigate how \textit{practical} such methods are - in practice, absolute convergence speed and total API costs depend on constants that aren't precisely foretold by our theoretical asymptotic analysis.

In addition, although these methods were initially developed for the case of decoding parameter perturbations, it would be interesting to investigate how robust these tests are when additional perturbations such as fine-tunes are introduced to the threat scenario. Even if the tests are no longer consistent, quantifying and bounding how confident the tests allow us to be is a potentially productive line of inquiry. 

Furthermore, we note that certain definitions and modeling decisions made in this paper are likely stronger than strictly necessary for the results to hold true - as a simple example, the definition of LLM equality required all model weights in LLM A and LLM B to be identical for them to be considered identical LLMs. However, the proofs in this paper relied on a constrained set of external model behaviors rather than the models' specific internal details. Thus, it may be meaningful to scrutinize how much we may weaken definitions and assumptions and still yield proofs of useful results.

Finally, although these methods are designed for LLM ownership identification in the stolen model threat scenario, which thus excludes deliberate impersonation attacks, it would be interesting to investigate how feasible impersonation attacks really are against the tests developed in this paper. As we argued in 5.5, the high level of redundancy in these tests creates a massive constraint satisfaction problem that seems, at surface level, highly difficult for an impersonation adversary to solve.

\subsection{Emerging trends: AI safety and security in agentic systems}
The framing of this paper was primarily that of IP theft prevention. Here, we briefly allude to the emerging practical scope of non-invasive black-box LLM ownership identification within the domain of AI safety. At its core, black-box LLM ownership verification is about providing 'proof' that an anonymous LLM is a model that it is pretending not to be. This maps usefully onto security problems that are receiving growing attention with the increasing adoption of agentic LLM systems - particularly the problem of 'Byzantine' agents, secret malicious models which have infiltrated agentic systems (\url{https://arxiv.org/pdf/2507.14928}, \url{https://dl.acm.org/doi/epdf/10.1145/3674399.3674445}, \url{https://arxiv.org/pdf/2408.00989v2}, \url{https://arxiv.org/pdf/2505.05103}). This new threat dimension calls for techniques that can flag and identify anonymous malicious models. As such, non-invasive, black-box techniques for LLM ownership identification, which do not assume that one has had an original hand in building the anonymous model, are a natural tool for tackling such problems. Refining rigorous understanding of which techniques work and when is thus becoming more important along a newfound axis, and invites further collaboration between safety, security, and provenance researchers.

\section{References}

TODO: Format citations

\section{Appendix}
\subsection{Enforcing $\dagger,2$}

The math behind the theorems in our "Theoretical Consequences of LLM A = LLM B" section hinges on $\displaystyle \dagger ,2$. In other words, we required that given any generated output string from LLM A, $\displaystyle t\in \{a_{1} ,...,a_{k} ,b_{1} ,...,b_{k}\}$ is a prefix of that string iff it is the first generated token of that string. $\displaystyle \dagger ,2$ is what makes the following 2 sets of definitions equivalent when LLM A = LLM B:



Version 1:
\begin{itemize}
\item Let $\displaystyle p_{j}( a,i)$ be the true probability of $\displaystyle a_{i}$ being a prefix string of a generated output from LLM B on context $\displaystyle C_{j}$ (unknown to us)
\item Let $\displaystyle p_{j}( b,i)$ be the true probability of $\displaystyle b_{i}$ being a prefix string of a generated output from LLM B on context $\displaystyle C_{j}$ (unknown to us)
\item Let $\displaystyle n_{j}( a,i)$ be the actual number of times $\displaystyle a_{i}$ appeared as a prefix string of a generated output from LLM B on context $\displaystyle C_{j}$ (known to us)
\item Let $\displaystyle n_{j}( b,i)$ be the actual number of times $\displaystyle b_{i}$ appeared as a prefix string of a generated output from LLM B on context $\displaystyle C_{j}$ (known to us)
\end{itemize}

Version 2:
\begin{itemize}
\item Let $\displaystyle p_{j}( a,i)$ be the true probability of $\displaystyle a_{i}$ in the first-token distribution of LLM B with respect to context $\displaystyle C_{j}$ (unknown to us)
\item Let $\displaystyle p_{j}( b,i)$ be the true probability of $\displaystyle b_{i}$ in the first-token distribution of LLM B with respect to context $\displaystyle C_{j}$ (unknown to us)
\item Let $\displaystyle n_{j}( a,i)$ be the actual number of times $\displaystyle a_{i}$ appeared as the first token of LLM B on context $\displaystyle C_{j}$ (known to us)
\item Let $\displaystyle n_{j}( b,i)$ be the actual number of times $\displaystyle b_{i}$ appeared as the first token of LLM B on context $\displaystyle C_{j}$ (known to us)
\end{itemize}



As an easy-to-understand motivating example, consider this situation: Suppose tokens '1', '2', and '12' are all in the vocabulary of LLM A . When the LLM returns string "12", we don't know if it returned '1' + '2' or '12'. What should we do to prevent this situation?



We will now discuss some practical ways to prevent this ambiguity. The first method works even when we don't have access to LLM A's vocabulary list.



\subsubsection{Method \#1: Forcing the LLM to output only 1 token}

LLM deployments often expose a max\_tokens parameter, putting a hard cap on the maximum number of tokens the LLM can generate before termination. The easiest way to prevent first token ambiguity is to set max\_tokens = 1.



In addition, many consumer-facing LLM providers often stream their output one at a time in string chunks that correspond to individual tokens. From this, we can, in principle, simulate max\_tokens = 1 by extracting the string representation of only the first token generated in the LLM's output string.



\subsubsection{Method \#2: Atomic Tokens}

Choosing $\displaystyle a_{i}$'s and $\displaystyle b_{i}$'s to be "atomic tokens" of the LLM is a non-"hacky" method to ensure that $\displaystyle a_{i}$/$\displaystyle b_{i}$ is the prefix of an output iff it is the first token of that output. Here, we define "atomic tokens" as tokens whose string representations are not a prefix of any other token's string representation and also do not contain any other token's string representation as a prefix of itself.



If we choose the $\displaystyle a_{i}$'s and $\displaystyle b_{i}$'s to be atomic tokens of the LLM, then whenever we see one of them appear as a prefix string of a generated output of the LLM, then we know with certainty that it is the first token of that string with zero ambiguity.



Since we have white-box access to LLM A (and therefore its vocabulary), we can in principle automate the process to find atomic tokens of LLM A.



Many LLMs share atomic tokens in common. Rare, standalone Unicode characters are often atomic tokens. For example:
\begin{itemize}
\item Certain emojis: Pushpin (U+1F4CC), Key (U+1F511), DNA (U+1F9EC), etc.
\item Certain symbols: Reference Mark (U+203B), Asterism (U+2042), Section Sign (U+00A7), etc.
\end{itemize}

\end{document}

